\documentclass[11pt, a4paper]{article}

%% == Packages =================================================================
\usepackage[utf8]{inputenc}
\usepackage[T1]{fontenc}

% Mathematics
\usepackage{amsmath, amssymb, amsthm, mathtools}
\usepackage{bm}

% Layout
\usepackage[margin=2.5cm]{geometry}
\usepackage{parskip}
\usepackage{enumitem}
\usepackage{booktabs}

% References
\usepackage{xcolor}
\usepackage[colorlinks=true, linkcolor=blue!60!black, citecolor=green!50!black,
            urlcolor=blue!70!black]{hyperref}
\usepackage{cleveref}

%% == Theorem environments =====================================================
\theoremstyle{definition}
\newtheorem{definition}{Definition}[section]
\newtheorem{assumption}{Assumption}[section]
\newtheorem{simplification}{Simplification}
\newtheorem{convention}{Convention}[section]

\theoremstyle{plain}
\newtheorem{theorem}{Theorem}[section]
\newtheorem{lemma}[theorem]{Lemma}
\newtheorem{proposition}[theorem]{Proposition}
\newtheorem{corollary}[theorem]{Corollary}

\theoremstyle{remark}
\newtheorem{remark}{Remark}[section]
\newtheorem{example}{Example}[section]

\crefname{assumption}{Assumption}{Assumptions}
\crefname{simplification}{Simplification}{Simplifications}
\crefname{convention}{Convention}{Conventions}
\crefname{remark}{Remark}{Remarks}
\crefname{definition}{Definition}{Definitions}
\crefname{example}{Example}{Examples}

%% == Notation shortcuts =======================================================
\newcommand{\R}{\mathbb{R}}
\newcommand{\E}{\mathbb{E}}
\newcommand{\Prob}{\mathbb{P}}
\newcommand{\Var}{\operatorname{Var}}
\newcommand{\Cov}{\operatorname{Cov}}
\newcommand{\tr}{\operatorname{tr}}
\newcommand{\dd}{\,\mathrm{d}}
\newcommand{\ind}[1]{\mathbf{1}\!\left[#1\right]}
\newcommand{\pos}[1]{\left(#1\right)^{\!+}}
\newcommand{\Lbar}{\bar{L}}
\newcommand{\Ltil}[1]{\tilde{L}_{#1}}

%% == Title ====================================================================
\title{Joint Latency Control for Synchronized Multi-Receiver Audio:\\
  Model, Loss, and the Structure of the Optimal Controller}
\author{}
\date{\today}

%% =============================================================================
\begin{document}
\maketitle

\begin{abstract}
This document states and solves a control problem for synchronized audio
streaming to $N$ receivers over a local network.  The problem has one loss
with four terms: mean latency (linear), cross-receiver synchronization error
(quadratic), buffer underrun (an indicator), and warp-rate roughness
(quadratic).  The disturbance is one stochastic process per receiver: the
per-sample delivery delay, modeled as baseline wander plus stationary
diffusion plus a pause process.  The document derives the structure of the
optimal controller.  The derivation is explicit about every simplification.
Each simplification gets a statement, a reason, a quantified defect, and a
telemetry monitor.  The result: the optimal controller factors into
per-receiver linear loops (a specific second-order law with damping ratio
$1/\sqrt2$) and a slow scalar program that positions the common latency
level.  The latency target is not an input of the problem.  It appears in
the solution as a dynamic state variable of the controller.
\end{abstract}

\tableofcontents
\newpage

%% =============================================================================
\section{Introduction}\label{sec:intro}

\subsection{The system in one page}\label{sec:system-one-page}

One sender streams audio over a local network to $N$ receivers.  Each
receiver drives one loudspeaker.  The sender packetizes samples and sends
them with their stream timestamps.  Each receiver holds arrived samples in
a buffer and plays them out through a resampler.  The resampler can play
slightly faster or slower than nominal rate.  This rate adjustment is
called \emph{warp}.  Warp is the only actuator in the system.

Three things can go wrong.
\begin{enumerate}[label=(\alph*)]
  \item \textbf{Latency.}  Audio plays late relative to capture.  For
    interactive sources (video, instruments) this is a direct cost.
  \item \textbf{Desynchronization.}  Two receivers play the same sample at
    different times.  The listener hears smeared or displaced sound.
  \item \textbf{Underrun.}  A receiver needs a sample that has not arrived.
    It must insert silence.  This is an audible dropout.
\end{enumerate}
A fourth cost is internal: warp changes pitch.  Fast warp changes are
audible.  Slow warp offsets are not.

These four costs pull in different directions.  Lower latency means a
smaller buffer, which means more underruns.  Tighter synchronization means
stiffer control, which means rougher warp.  The purpose of this document is
to treat the four costs as one optimization problem and to derive the
controller that this problem implies.

\subsection{The three goals are one problem}\label{sec:one-problem}

Historically the three goals were handled by three mechanisms: a fixed
latency target chosen by hand (latency), independent per-receiver
controllers holding that target (synchronization), and a safety margin
inside the target (underrun).  This document shows that the three
mechanisms are the correct \emph{shape} of the solution, but that the
target must not be an input.  The optimization produces the target as an
output.  The target becomes an internal, slowly moving state of the
controller, computed from measured statistics of the network.

\subsection{How to read this document}\label{sec:how-to-read}

Sections \ref{sec:model}--\ref{sec:loss} state the problem: the physical
model, the stochastic process, and the loss.  These sections contain no
approximations beyond three explicitly flagged identifications: an
$O(10^{-4})$ actuator rescaling (\cref{prop:latency-dynamics}), the
margin--queue identification to one packet period
(\cref{rem:margin-queue}), and the zero-drain ramp idealization
(\cref{prop:ramp}).  Section \ref{sec:sufficiency} proves which statistics
of the process each loss term can see.  This section explains \emph{why}
the solution splits into two different kinds of mathematics.
Section \ref{sec:reductions} performs exact coordinate changes.
Section \ref{sec:quadratic} solves the quadratic part of the problem with
linear-quadratic machinery.  Section \ref{sec:indicator} handles the
underrun term.  Every simplification lives in a numbered environment
(\textbf{S1}, \textbf{S2}, \dots) with a quantified defect.
Section \ref{sec:composition} composes the two solutions and proves the
decomposition result.  Section \ref{sec:controller} states the resulting
controller as an algorithm.  Section \ref{sec:register} collects all
simplifications in one table.  Section \ref{sec:limits} lists what the
derivation does \emph{not} cover.

Two style rules apply throughout.  First: every symbol is defined before
use, and \cref{sec:notation} collects the recurring ones.  Second: every claim is
either proven here, or imported with an explicit citation and marked as
imported.  There are no silent steps.

%% =============================================================================
\section{The System Model}\label{sec:model}

\subsection{Timelines, samples, and schedules}\label{sec:timelines}

\begin{convention}[Units]\label{conv:units}
All times are in seconds.  Stream position is measured in seconds of audio,
not in sample counts.  This removes the sample rate $F_s$ from all
formulas.  To convert to sample units, multiply seconds of audio by $F_s$.
\end{convention}

\begin{definition}[Stream and emission]\label{def:stream}
The audio stream is indexed by \emph{stream time} $s \in \R$: the position
inside the audio signal, in seconds of audio.  The sender emits (captures
and transmits) stream position $s$ at wall time $\tau(s)$.  For an ideal
sender clock, $\tau(s) = s + \text{const}$; we absorb the constant and
write $\tau(s) = s + c_A(s)$, where $c_A$ is the sender clock error, a
slowly varying function with $|\dot c_A| \le 10^{-4}$ (clock rates differ
from nominal by at most tens of ppm).
\end{definition}

\begin{definition}[Arrival]\label{def:arrival}
Stream position $s$ becomes available in receiver $i$'s buffer at wall
time
\begin{equation}\label{eq:arrival}
  a_i(s) \;=\; \tau(s) + D_i(s),
\end{equation}
where $D_i(s) \ge 0$ is the \emph{delivery delay} of receiver $i$ for
stream position $s$.  The delay $D_i$ covers everything between emission
and availability: packetization, network transit, queueing, and any
receiver-side processing before the sample is playable.  The delay
processes $(D_1, \dots, D_N)$ are the stochastic disturbance of the
problem.  \Cref{sec:process} models them.
\end{definition}

\begin{definition}[Playout schedule and warp]\label{def:schedule}
Receiver $i$ chooses a strictly increasing \emph{playout schedule}
$p_i(s)$: stream position $s$ plays at wall time $p_i(s)$.  Let $s_i(t)$
denote the inverse: the stream position scheduled at wall time~$t$, so
$p_i(s_i(t)) = t$.  The \emph{warp} is the fractional playout rate offset
\begin{equation}\label{eq:warp-def}
  u_i(t) \;\coloneqq\; \frac{\dd s_i}{\dd t}(t)\,\bigl(1 + \dot c_{B,i}\bigr) - 1 ,
\end{equation}
where $c_{B,i}$ is receiver $i$'s clock error function, defined like
$c_A$ in \cref{def:stream}, with the same ppm-scale rate bound.  In words: $u_i > 0$ means receiver $i$ consumes the stream faster
than nominal; $u_i < 0$ means slower.  The warp is bounded by the
resampler: $|u_i| \le u_{\max}$, with $u_{\max} = 5\times 10^{-3}$ in the
implementation.
\end{definition}

\begin{definition}[End-to-end latency]\label{def:latency}
The \emph{latency} of receiver $i$ at wall time $t$ is the age of the
sample it is scheduled to play:
\begin{equation}\label{eq:latency-def}
  L_i(t) \;\coloneqq\; t - \tau\!\bigl(s_i(t)\bigr).
\end{equation}
\end{definition}

\begin{proposition}[Latency dynamics]\label{prop:latency-dynamics}
The latency obeys
\begin{equation}\label{eq:latency-dynamics}
  \dot L_i(t) \;=\; w_i(t) - u_i(t),
  \qquad
  |w_i(t)| \;\le\; 3\times 10^{-4},
\end{equation}
where $w_i \coloneqq
(\dot c_{B,i} - \dot c_A)/(1 + \dot c_{B,i})$ is the relative clock
drift between receiver $i$ and the sender.
\end{proposition}

\begin{proof}
Differentiate \cref{eq:latency-def}:
$\dot L_i = 1 - \tau'(s_i(t))\,\dot s_i(t)
          = 1 - \bigl(1+\dot c_A\bigr)\dot s_i$.
By \cref{eq:warp-def}, $\dot s_i = (1+u_i)/(1+\dot c_{B,i})$.  Therefore
\[
  \dot L_i
  = 1 - \frac{(1+\dot c_A)(1+u_i)}{1+\dot c_{B,i}}
  = \frac{\dot c_{B,i} - \dot c_A - u_i\,(1 + \dot c_A)}{1 + \dot c_{B,i}} .
\]
Define $w_i \coloneqq (\dot c_{B,i}-\dot c_A)/(1+\dot c_{B,i})$ and absorb
the factor $(1+\dot c_A)/(1+\dot c_{B,i}) = 1 + O(10^{-4})$ into the
definition of $u_i$ (a relative rescaling of the actuator by at most
$10^{-4}$, far below actuator resolution).  Both clock rates are within
$10^{-4}$ of nominal, so
$|w_i| \le 2\times 10^{-4}/(1 - 10^{-4}) < 3\times 10^{-4}$.
\end{proof}

\begin{remark}[What \cref{eq:latency-dynamics} says]
Latency is an integrator driven by clock drift and steered by warp.  The
controller cannot set latency directly.  It can only slide latency around
at up to $u_{\max} = 5$\,ms per second.  This bound is the single most
important quantitative fact in this document.  It reappears in
Simplification~S3, where it decouples the underrun term from the fast
control loop, and in \cref{sec:glide}, where it limits how fast the
latency level can be raised.
\end{remark}

\subsection{The buffer and underrun}\label{sec:buffer}

\begin{definition}[Underrun]\label{def:underrun}
Receiver $i$ is \emph{in underrun} at wall time $t$ when the sample it is
scheduled to play has not arrived:
\begin{equation}\label{eq:underrun-def}
  \text{underrun}_i(t)
  \;\iff\;
  a_i\!\bigl(s_i(t)\bigr) \;>\; t .
\end{equation}
\end{definition}

\begin{convention}[Silence convention]\label{conv:silence}
During underrun the schedule keeps advancing and the receiver outputs
silence for the missing samples.  The schedule is not paused and not
reset.  Missing audio is counted by the loss (\cref{sec:loss}); it is not
compensated later.
\end{convention}

\begin{proposition}[Underrun is a delay-versus-latency comparison]
\label{prop:underrun-margin}
For all $i, t$:
\begin{equation}\label{eq:underrun-margin}
  \text{underrun}_i(t)
  \;\iff\;
  D_i\!\bigl(s_i(t)\bigr) \;>\; L_i(t).
\end{equation}
\end{proposition}

\begin{proof}
By \cref{eq:arrival}, $a_i(s_i(t)) > t$ iff
$\tau(s_i(t)) + D_i(s_i(t)) > t$ iff
$D_i(s_i(t)) > t - \tau(s_i(t)) = L_i(t)$ by \cref{eq:latency-def}.
\end{proof}

\begin{definition}[Margin]\label{def:margin}
The \emph{margin} of receiver $i$ at wall time $t$ is
\begin{equation}\label{eq:margin-def}
  m_i(t) \;\coloneqq\; L_i(t) - D_i\!\bigl(s_i(t)\bigl).
\end{equation}
Underrun at $t$ is exactly $m_i(t) < 0$.
\end{definition}

\begin{remark}[Margin and buffer depth]\label{rem:margin-queue}
When arrivals are in stream order, the buffer of receiver $i$ at time $t$
holds exactly the samples $s$ with $s_i(t) \le s$ and $a_i(s) \le t$, and
the buffered duration equals the margin up to the local variation of $D_i$
over one buffer length.  Packet reordering breaks exact equality but the
identification ``margin $\approx$ buffered audio duration'' is accurate to
one packet period.  We use the word \emph{margin} for the model quantity
and \emph{queue depth} for the implementation measurement.
\end{remark}

\subsection{Measurement: two feedback architectures}\label{sec:measurement}

The controller observes the system through one of two measured outputs.
Both exist in the implementation; the choice is per stream.

\begin{definition}[Latency measurement, ``e2e'' backend]\label{def:meas-e2e}
The receiver measures
\begin{equation}\label{eq:meas-e2e}
  y_i(t) \;=\; L_i(t) + n_i(t),
\end{equation}
where $n_i$ is measurement error.  Its dominant component is the error of
the clock synchronization between sender and receiver: computing
$t - \tau(s)$ requires comparing the two clocks.  On hosts disciplined by
the same time service, $n_i$ is microseconds.  On an undisciplined host,
$n_i$ can drift by milliseconds over minutes; this drift enters the loop
as a false disturbance.  \Cref{sec:transfer} quantifies the consequence.
\end{definition}

\begin{definition}[Margin measurement, ``niq'' backend]\label{def:meas-niq}
The receiver measures its own buffered duration:
\begin{equation}\label{eq:meas-niq}
  y_i(t) \;=\; m_i(t).
\end{equation}
This measurement needs no cross-host clock comparison.  It is exact up to
packet quantization.  Its cost: the measured signal now contains the
delay process, so network disturbances drive the control loop
(\cref{sec:transfer}).
\end{definition}

%% =============================================================================
\section{The Delay Process}\label{sec:process}

This section models the disturbance $(D_1, \dots, D_N)$.  The model is the
result of measurement, not of convenience; each component corresponds to a
distinct physical mechanism, and the telemetry estimates each component
separately.

\subsection{The three components}\label{sec:components}

\begin{assumption}[Additive decomposition]\label{as:decomposition}
For each receiver $i$,
\begin{equation}\label{eq:decomposition}
  D_i(s) \;=\; B_i(t_s) \;+\; \varepsilon_i(s) \;+\; P_i(s),
  \qquad t_s = \tau(s),
\end{equation}
where $B_i$ is the \emph{baseline}, $\varepsilon_i$ is the
\emph{diffusion}, and $P_i$ is the \emph{pause component}, defined in
\cref{as:baseline,as:diffusion,as:pause} below.
\end{assumption}

\begin{assumption}[Baseline]\label{as:baseline}
$B_i(t)$ is differentiable and slow: its derivative is bounded,
$|\dot B_i| \le \bar\beta$ with $\bar\beta$ of order $10^{-5}$ (tens of
microseconds per second), and its power spectrum is concentrated below the
control bandwidth $\kappa$ defined in \cref{sec:lqr}.  Physically:
propagation and processing floor, thermal drift of clocks and interface
timing, slow route or load changes.
\end{assumption}

\begin{assumption}[Diffusion]\label{as:diffusion}
$\varepsilon_i(s)$ is stationary, mean-zero, and mixing (correlations
decay; there is a horizon beyond which values are almost independent).
It has variance $\sigma_i^2$ and power spectral density $S_{\varepsilon,i}(\omega)$.
Its tail is light: there are constants $x_0, C$ such that
$\Prob(\varepsilon_i > x) \le C e^{-x/x_0}$ for all $x>0$.  Physically:
interrupt coalescing, scheduler granularity, cross-traffic queueing at
sub-millisecond scale.  Measured values in the reference deployment:
$\sigma_i$ from tens of microseconds up to about one millisecond on wired
hosts.  The upper end matters: at an aggressive latency level of order
$10$\,ms the margin is only a few multiples of $\sigma_i$, so the tail of
$\varepsilon_i$ is a first-class contribution to underrun risk, not a
negligible one (\cref{sec:risk-evaluation}).
\end{assumption}

\begin{assumption}[Pause component]\label{as:pause}
Delivery to receiver $i$ stops entirely during \emph{pause episodes}.
Episode $k$ starts at wall time $T_{i,k}$ and lasts $H_{i,k}$ seconds
(the \emph{height}); delivery resumes at $T_{i,k} + H_{i,k}$.  Stream
content emitted during the episode is delivered late (after resume),
recovered by forward error correction, or lost; the split between these
outcomes is a property of the episode mechanism.  The pairs
$(T_{i,k}, H_{i,k})$ form a stationary marked point process with
intensity $\lambda_i$ (episodes per second) and mark law $F_i$ (the Palm
distribution of heights; $\bar F_i \coloneqq 1 - F_i$).  Clustering of
episode times is allowed and is a measured property, not a parameter.
Physically: radio retraining, roaming, bulk interference, CPU stalls.
Measured heights range from milliseconds (wired hosts) to $10^2$--$10^3$
milliseconds (radio hosts).
\end{assumption}

\begin{proposition}[Pause ramp shape]\label{prop:ramp}
Under \cref{as:pause} and the zero-drain idealization below, during
episode $k$ the pause component is a descending ramp: for stream
positions $s$ with $\tau(s) \in [T_{i,k},\, T_{i,k}+H_{i,k}]$,
\begin{equation}\label{eq:ramp}
  P_i(s) \;=\; T_{i,k} + H_{i,k} - \tau(s),
\end{equation}
and $P_i(s) = 0$ outside episodes.  In particular
$\max_s P_i(s) = H_{i,k}$ during episode $k$: the height equals both the
duration of the episode and the peak extra delay.

\emph{Zero-drain idealization:} withheld content is treated as arriving
at the instant delivery resumes.  In reality the burst drains over the
time of \cref{as:recovery}(a) (milliseconds on a local network); the
ramp is exact up to that drain time.  This idealization is carried as
part of Simplification S2's defect accounting
(\cref{sec:simplification-chain}).
\end{proposition}

\begin{proof}
A sample emitted at $\tau(s)$ inside the episode cannot arrive before
delivery resumes at $T_{i,k}+H_{i,k}$.  If it is delivered at all, it
arrives in the burst of withheld content just after resume --- at the
resume instant, under the idealization --- so its extra delay is
$T_{i,k}+H_{i,k}-\tau(s)$, which falls linearly from $H_{i,k}$ (for the
sample emitted at episode start) to $0$ (at episode end).  Samples
emitted outside episodes have no withheld path, so $P_i = 0$.
\end{proof}

\begin{definition}[Recovery time]\label{def:recovery}
The \emph{recovery time} $\rho_{i,k}$ of episode $k$ is the time from
delivery resume until the margin of receiver $i$ is back at its
pre-episode value.
\end{definition}

\begin{assumption}[Recovery mechanisms]\label{as:recovery}
No instant recovery is assumed.  Two mechanisms bound $\rho_{i,k}$:
\begin{enumerate}[label=(\alph*)]
  \item \textbf{Redelivery burst (fast).}  When the withheld content was
    queued upstream, it arrives in a burst at the link rate once delivery
    resumes.  The burst refills the buffer directly:
    $\rho_{i,k} \approx H_{i,k}\cdot R_{\mathrm{stream}}/R_{\mathrm{link}}$
    (stream bit rate over link capacity), milliseconds on a local
    network.  This is the measured behavior of
    radio dropouts in the reference deployment: the instantaneous queue
    depth spikes by approximately the episode height right after resume.
  \item \textbf{Warp-limited restoration (slow).}  When content was lost,
    or when the playout schedule slipped during the episode (a stalled
    playout, or a controller that converted the transient into latency),
    the deficit can only be worked off by warp:
    $\rho_{i,k} \le \chi_{i,k}/u_{\max}$, where
    $\chi_{i,k} = \min(H_{i,k}, m_i)$ is the margin the episode consumed.
    With $u_{\max} = 5\times10^{-3}$, a consumed margin of $10$\,ms takes
    up to $2$\,s to restore; $100$\,ms takes up to $20$\,s.
\end{enumerate}
Which mechanism applies is a property of the episode and of the
implementation path; $\rho$ is treated as a measured per-episode mark.
The interval after an episode during which the margin is still depressed
is a window of raised vulnerability; \cref{sec:risk-evaluation} prices a
following episode inside this window using the measured gap law.
\end{assumption}

\begin{assumption}[Cross-receiver structure]\label{as:cross}
The joint process $(D_1,\dots,D_N)$ is stationary.  Components decompose
along the network path tree: each of $B, \varepsilon, P$ splits into a
part common to all receivers (sender host, shared switch), parts shared by
subtrees, and idiosyncratic parts.  No independence across receivers is
assumed anywhere.  Measured fact: pause episodes co-occur across receivers
at rates two orders of magnitude above the independence prediction.
\end{assumption}

\subsection{What the model does not assume}\label{sec:not-assumed}

The model does not assume Poisson pause times (clustering is measured).
It does not assume any specific height law $F_i$ (heavy or multimodal laws
are allowed; heights are bounded in practice by session supervision, so
all moments of $D_i$ exist).  It does not assume Gaussian diffusion (only
light tails and mixing).  It does not assume independence across
receivers.  It does not assume instant recovery after an episode
(\cref{as:recovery}), and it does not assume that the diffusion tail is
negligible at the operating margin.

\subsection{Levels, not increments}\label{sec:levels-not-increments}

Classical jitter metrics are built from \emph{increments}: differences of
consecutive delays.  The following elementary fact shows that increments
cannot drive a buffer model, and it calibrates what the white-noise terms
in \cref{sec:quadratic} may and may not represent.

\begin{proposition}[Increment sums telescope]\label{prop:telescope}
Let $\eta(n)$, $n \in \mathbb{Z}$, be any stationary sequence (per-packet
delay levels) with variance $\sigma_\eta^2 < \infty$, and let
$\xi(n) = \eta(n)-\eta(n-1)$.  Then for every $M$,
\[
  \sum_{n=1}^{M} \xi(n) \;=\; \eta(M) - \eta(0),
  \qquad
  \Var\!\Bigl[\sum_{n=1}^{M}\xi(n)\Bigr] \;\le\; 4\sigma_\eta^2 ,
\]
so the long-run variance rate of the partial sums is zero:
$\lim_{M\to\infty} \Var[\sum_1^M \xi]/M = 0$.
\end{proposition}

\begin{proof}
The sum telescopes term by term.  The variance of a difference of two
variables, each with variance $\sigma_\eta^2$, is at most
$4\sigma_\eta^2$ by the Cauchy--Schwarz inequality.  Dividing a bounded
quantity by $M$ gives zero in the limit.
\end{proof}

\begin{corollary}\label{cor:no-brownian}
A stationary delay process contributes no Brownian (accumulating) noise to
the buffer.  The buffer responds to the delay \emph{level} directly:
a delay change of $x$ seconds moves the margin by $-x$ seconds, once, and
the margin recovers when the delay recovers.  Any model that drives the
buffer with white noise of nonzero intensity implicitly assumes a
\emph{nonstationary} (wandering) delay level.  In this document, white
noise terms therefore represent only the baseline wander $B_i$ and
measurement noise; the stationary diffusion $\varepsilon_i$ enters through
its spectrum and the pause component $P_i$ through the underrun term.
\end{corollary}

%% =============================================================================
\section{The Loss}\label{sec:loss}

\subsection{Policies}\label{sec:policies}

\begin{definition}[Causal policy]\label{def:policy}
A \emph{policy} $\pi$ is a rule that chooses every receiver's warp
$u_i(t)$ as a function of observations available at time $t$: the
receiver's own measured output $y_i$ up to $t$
(\cref{def:meas-e2e,def:meas-niq}), its own arrival history, and any
messages received from the sender by time $t$.  Messages from receivers reach the sender after a one-way network
delay, and sender broadcasts reach receivers after another; any
closed information loop through the sender takes at least a round
trip.  The information structure is therefore
distributed: no receiver sees another receiver's state directly, and the
sender sees receiver reports with delay.
\end{definition}

\subsection{The loss functional}\label{sec:loss-functional}

Let $\omega$ denote one realization of the joint disturbance: the delay
processes of all receivers and the clock drifts.  All randomness lives in
$\omega$.

\begin{definition}[The loss]\label{def:loss}
For a policy $\pi$, define
\begin{equation}\label{eq:loss}
  J(\pi) \;=\;
  \E_\omega\!\left[\;
    \limsup_{T\to\infty}\,
    \frac{1}{T}\int_0^{T} \Lambda(t,\omega)\;\dd t
  \right],
\end{equation}
with the instantaneous loss rate, evaluated on the realization,
\begin{equation}\label{eq:pointwise-loss}
  \Lambda(t,\omega)
  \;=\;
  \alpha\,\Lbar(t)
  \;+\;
  \frac{q}{N}\sum_{i=1}^{N}\bigl(L_i(t)-\Lbar(t)\bigr)^2
  \;+\;
  \frac{c}{N}\sum_{i=1}^{N}\ind{m_i(t) < 0}
  \;+\;
  \frac{r}{N}\sum_{i=1}^{N}\dot u_i(t)^2 ,
\end{equation}
where $\Lbar(t) = \frac1N\sum_i L_i(t)$ is the instantaneous fleet mean
latency and $\alpha, q, c, r > 0$ are weights.
\end{definition}

The four terms, in words:
\begin{enumerate}[label=(\roman*)]
  \item \textbf{Latency (linear).}  Every second of average latency costs
    $\alpha$ per second of operation.  Latency is bad in proportion to its
    size; there is no threshold and no target.
  \item \textbf{Synchronization (quadratic).}  The instantaneous variance
    of latency across receivers, on this realization.  This is the
    physical desynchronization the listener hears.  It contains no
    reference value: it compares receivers with each other only.
  \item \textbf{Underrun (indicator).}  The fraction of receivers
    currently playing silence for lack of data.  By
    \cref{prop:underrun-margin} this is a property of delay versus
    latency; by \cref{conv:silence} silence duration is what accumulates.
  \item \textbf{Warp roughness (quadratic).}  Fast warp changes are
    audible pitch transients.  A constant warp is free: only the
    derivative is charged.  This choice has a structural consequence used
    in \cref{sec:composition}: the loss puts zero price on any constant
    warp offset, hence zero price on holding any fixed latency level.
\end{enumerate}

\begin{remark}[Units]\label{rem:units}
Let $\mathsf{u}$ denote one loss unit.  Then $\Lambda$ has units
$\mathsf{u}/\mathrm{s}$ and
\[
  [\alpha] = \mathsf{u}\,\mathrm{s}^{-2},\quad
  [q] = \mathsf{u}\,\mathrm{s}^{-3},\quad
  [c] = \mathsf{u}\,\mathrm{s}^{-1},\quad
  [r] = \mathsf{u}\,\mathrm{s}.
\]
Two combinations will matter:
$\alpha/c$ has units $\mathrm{s}^{-1}$ (a rate --- it will equal the
deviation's occupation density at the optimum,
\cref{sec:level-program}; for the pause component this is an episode
exceedance rate), and $(q/r)^{1/4}$ has units
$\mathrm{s}^{-1}$ (a bandwidth --- it will be the fast loop's natural
frequency in \cref{sec:lqr}).  Dimensional consistency of both is a
useful check on the derivation.
\end{remark}

\subsection{Modeling choices, stated openly}\label{sec:modeling-choices}

Three choices in \cref{def:loss} are choices, not necessities.

\begin{enumerate}[label=(\alph*)]
  \item \textbf{Mean over receivers in the latency term.}  The maximum
    over receivers is an alternative.  With the synchronization term
    active, receivers sit within a fraction of a millisecond of each
    other, so mean and maximum differ by less than the resolution of any
    latency preference.  The mean gives cleaner algebra.
  \item \textbf{Linear latency cost.}  Any increasing function of
    latency is admissible.  Linearity makes $\alpha$ interpretable (the
    price of one second of latency) and, in \cref{sec:level-program},
    makes the operating point depend on the single ratio $\alpha/c$.
  \item \textbf{Underrun counted by time, symmetrically across
    receivers.}  The indicator charges silence duration.  It does not
    separately charge the \emph{number} of dropouts, and it does not
    model listener adaptation.  A count-based or risk-averse variant
    changes \cref{sec:indicator}; \cref{sec:limits} discusses where.
\end{enumerate}

\begin{remark}[What is absent: session restarts]\label{rem:no-restart}
The implementation supervises sessions and restarts them on gross
failures (unbounded queue growth, tolerance violations).  A restart is a
loud, discrete event that this loss does not model beyond the silence it
causes.  The derivation therefore optimizes normal operation; supervision
remains a separate safety layer.  See \cref{sec:limits}.
\end{remark}

%% =============================================================================
\section{What Each Loss Term Can See}\label{sec:sufficiency}

The four loss terms differ in which statistics of the disturbance they
depend on.  This section makes that precise.  The result explains the
architecture of the whole solution: three terms live inside a closed
``moment algebra'' that linear-quadratic theory solves, and one term ---
the indicator --- provably lives outside it.

\subsection{Linear and quadratic terms need only moments}
\label{sec:moments-suffice}

\begin{proposition}[First moments suffice for linear terms]
\label{prop:linear-first-moment}
Let $x$ be a random vector with mean $\mu$, and let $g(x) = a^\top x$ be
linear.  Then $\E[g(x)] = a^\top\mu$ depends on the law of $x$ only
through its first moment.
\end{proposition}

\begin{proof} Linearity of expectation. \end{proof}

\begin{proposition}[Second moments suffice for quadratic terms]
\label{prop:quadratic-second-moment}
Let $x$ be a random vector with mean $\mu$ and covariance $\Sigma$, and
let $g(x) = x^\top Q x$.  Then
\begin{equation}\label{eq:trace-formula}
  \E[g(x)] \;=\; \tr(Q\Sigma) + \mu^\top Q \mu ,
\end{equation}
which depends on the law of $x$ only through its first two moments.
\end{proposition}

\begin{proof}
Write $x = \mu + z$ with $\E[z]=0$, $\E[zz^\top]=\Sigma$.  Then
$\E[x^\top Qx] = \mu^\top Q\mu + 2\mu^\top Q\,\E[z] + \E[z^\top Qz]$,
and $\E[z^\top Qz] = \E[\tr(Qzz^\top)] = \tr(Q\Sigma)$ by linearity of
the trace and of expectation.
\end{proof}

For stationary processes passed through linear time-invariant systems,
``second moments'' means the power spectral density; \cref{sec:spectral}
gives the exact statement.  The consequence: for the latency,
synchronization, and warp terms, two disturbance processes with equal
first and second moments produce \emph{identical} loss under every linear
policy.  Everything else about the disturbance is invisible to these
terms.

\subsection{The indicator term is not determined by moments}
\label{sec:indicator-not-moments}

\begin{proposition}[Moment insufficiency, marginal form]
\label{prop:counterexample}
There exist two stationary processes with equal mean and equal
autocovariance function whose stationary underrun fractions at the same
margin differ by more than eighty orders of magnitude.
\end{proposition}

\begin{proof}
Take the pause process of \cref{as:pause} with constant intensity
$\lambda = 10^{-2}\,\mathrm{s}^{-1}$, deterministic height
$H \equiv 0.5$\,s, non-overlapping episodes, and no other components.
By \cref{prop:ramp}, during an episode the level descends linearly from
$H$ to $0$ at unit rate, so the fraction of time the level exceeds
$m = 0.4$\,s is, by \cref{prop:campbell} below,
$\lambda\,\E\pos{H - m} = 10^{-2}\times 0.1 = 10^{-3}$.
Its stationary mean is
$\mu_P = \lambda \int_0^H (H-x)\dd x = \lambda H^2/2 = 1.25\times10^{-3}$\,s,
and its stationary second moment is
$\lambda\int_0^H (H-x)^2\dd x = \lambda H^3/3 = 4.17\times10^{-4}\,\mathrm{s}^2$,
giving variance $\approx 4.15\times10^{-4}\,\mathrm{s}^2$, standard
deviation $\sigma_P \approx 20.4$\,ms.  Its autocovariance function is
some fixed function $C(\cdot)$.

Now take the Gaussian process with the same mean $\mu_P$ and the same
autocovariance $C$ (one exists for any valid autocovariance).  Its
stationary exceedance of $m=0.4$\,s is
$\bar\Phi\bigl((0.4-0.00125)/0.0204\bigr) = \bar\Phi(19.5) < 10^{-84}$,
using the Mills-ratio bound
$\bar\Phi(x) \le \varphi(x)/x = e^{-x^2/2}/(x\sqrt{2\pi})$.  The two processes agree on all first
and second moments and disagree on the underrun fraction by a factor
exceeding $10^{80}$.
\end{proof}

\begin{remark}
The example is extreme by design, but the mechanism is the measured one:
pause heights are hundreds of standard deviations of the bulk.  Any
moment-matched Gaussian surrogate concludes that underruns never happen.
\end{remark}

\subsection{Matching the marginal is still not enough for event structure}
\label{sec:copula}

A natural repair: transform a Gaussian process through a monotone map so
that the \emph{marginal} distribution matches exactly (a
``trans-Gaussian'' or Gaussian-copula model), and match the covariance as
well.  This fixes the exceedance \emph{fraction} by construction.  It
still fails on the \emph{time structure} of exceedances, and the failure
is qualitative.

\begin{proposition}[Sojourn shrinkage for Gaussian-based models]
\label{prop:sojourn}
Let $X$ be a stationary, almost-surely differentiable Gaussian process
with variance $\sigma_0^2$ and derivative variance $\sigma_1^2$.  The
mean duration of one visit of $X$ above level $z$ satisfies
\begin{equation}\label{eq:sojourn}
  \mu_{\mathrm{soj}}(z)
  \;=\;
  \frac{\Prob(X > z)}{N^+(z)}
  \;=\;
  \frac{\sqrt{2\pi}\,\sigma_0^2}{\sigma_1\, z}\,\bigl(1+o(1)\bigr)
  \;\xrightarrow[z\to\infty]{}\; 0,
\end{equation}
where $N^+(z)$ is the rate of upcrossings of level $z$.  A monotone
marginal transform changes the level scale but not the vanishing:
high-level visits of any trans-Gaussian process built on a smooth
(almost surely differentiable) Gaussian base become arbitrarily short.
By contrast, the pause process of \cref{as:pause} has mean visit duration
$\E[\,H - m \mid H > m\,]$ above margin $m$: a quantity of the order of
the height spread, bounded away from zero.
\end{proposition}

\begin{proof}
Two imports, both collected in \cref{sec:app-rice}: Rice's formula for
the upcrossing rate, and the ergodic occupation--crossing ratio
identity $\mu_{\mathrm{soj}}(z) = \Prob(X>z)/N^+(z)$ (mean visit
duration equals stationary exceedance probability divided by visit
rate).  Rice's formula gives
$N^+(z) = \frac{\sigma_1}{2\pi\sigma_0} e^{-z^2/2\sigma_0^2}$.
The Mills ratio bound
$\bar\Phi(x) = \frac{\varphi(x)}{x}(1+O(x^{-2}))$ with $x = z/\sigma_0$
gives
$\mu_{\mathrm{soj}}
 = \bar\Phi(x)\,\frac{2\pi\sigma_0}{\sigma_1}e^{x^2/2}
 = \frac{1}{\sqrt{2\pi}\,x}\,\frac{2\pi\sigma_0}{\sigma_1}\,(1+o(1))
 = \frac{\sqrt{2\pi}\,\sigma_0^2}{\sigma_1 z}(1+o(1))$.
The pause statement is \cref{prop:ramp}: one episode with height $H > m$
stays above $m$ for exactly $H - m$.  Underlying both is the
tail-independence property of the Gaussian copula
(\cref{sec:app-tail-independence}).
\end{proof}

\begin{remark}[Scope of the sojourn argument]\label{rem:sojourn-scope}
\Cref{prop:sojourn} is an asymptotic statement ($z \to \infty$).  At a
\emph{fixed} operating margin it supplies the mechanism and the
direction, not a bound: how short the transformed process's visits are
at a given level depends on the transform.  The operational claim ---
that marginal-matched Gaussian-based models mispredict event durations
at deployed margins --- rests on the exact fixed-level counterexample
(\cref{prop:counterexample}) for magnitudes and on the measured
event-duration histograms of the telemetry, with the asymptotics
pinning the reason.
\end{remark}

\subsection{Where the temporal law does and does not matter}
\label{sec:where-temporal}

The last two propositions justify modeling pauses explicitly.  The
following exactness result then draws a perhaps surprising boundary: for
the \emph{mean} underrun loss, only marginal pause statistics matter.

\begin{proposition}[Campbell evaluation of the underrun fraction]
\label{prop:campbell}
Under \cref{as:pause}, suppose episodes of receiver $i$ do not overlap,
the margin is held at the constant value $m$ throughout, and recovery is
instant ($\rho = 0$; \cref{sec:indicator} quantifies the error of all
three suppositions).  Then the long-run fraction of time in underrun
caused by pauses alone is
\begin{equation}\label{eq:campbell}
  R_i(m)
  \;=\;
  \lambda_i \,\E\pos{H_i - m},
\end{equation}
where the expectation is over the Palm mark law $F_i$.  This holds for
\emph{any} stationary marked point process: clustering of episode times
does not change the mean underrun fraction.
\end{proposition}

\begin{proof}
By \cref{prop:ramp} and \cref{def:margin}, during episode $k$ the margin
is $m - P_i$, which is negative exactly while $P_i > m$, i.e., for
duration $\pos{H_{i,k}-m}$ (the ramp descends at unit rate).  Episodes do
not overlap, so total underrun time in $[0,T]$ is
$\sum_{k: T_{i,k}\in[0,T]} \pos{H_{i,k}-m}$ up to edge effects of two
episodes.  The Campbell--Mecke formula for stationary marked point
processes (imported; statement and reference in \cref{sec:app-campbell})
gives the expectation of this sum as
$\lambda_i T\,\E\pos{H_i - m}$.  Dividing by $T$ and letting
$T \to \infty$ gives \cref{eq:campbell}; the limit holds almost surely
under ergodicity by the same theorem's ergodic form.
\end{proof}

\begin{remark}[The division of labor]\label{rem:division}
Combining \cref{prop:counterexample,prop:sojourn,prop:campbell}:
\begin{itemize}
  \item Under instant recovery, the \emph{mean} pause-caused underrun
    loss at a fixed margin needs only the marginal exceedance structure
    $(\lambda_i, F_i)$.  Clustering does not move it.
  \item With finite recovery times (\cref{as:recovery}) the mean loss
    gains a compound term: an episode arriving while the margin is still
    depressed does extra damage.  This term depends on the \emph{gap
    law} between consecutive episodes, so the temporal structure enters
    the mean loss through exactly one channel --- gaps shorter than
    recovery times.  \Cref{sec:risk-evaluation} computes it.
  \item Beyond the mean, the temporal law also governs the variance of
    realized loss over finite horizons (hence how fast
    $(\lambda_i, F_i)$ can be estimated and how wide the confidence
    bands of the level program are), the validity of Simplifications
    S3--S5, and any risk-averse variant of the loss.
\end{itemize}
This sharpens the reason for measuring gap structure: it prices the
compound risk, it calibrates the trust in the estimated risk, and it
validates the model.
\end{remark}

%% =============================================================================
\section{Exact Reductions}\label{sec:reductions}

This section performs coordinate changes that are exact, preceded by
the one statistical convention everything later rests on (S1 ---
stated first so that its scope is visible, and itself carrying a
defect like every other S-item).  The output is a decomposition of the loss into parts that
the quadratic machinery of \cref{sec:quadratic} can process, plus a
precise identification of the one term that resists.

\subsection{From time averages to stationary expectations}
\label{sec:ergodic}

\begin{simplification}[S1: stationary ergodic regime]\label{simp:ergodic}
\textbf{Statement.}  Within one operating regime, the closed-loop system
under a fixed policy is stationary and ergodic: the time average in
\cref{eq:loss} equals the expectation of $\Lambda$ under the stationary law,
\[
  J(\pi) \;=\; \E\bigl[\Lambda(0,\omega)\bigr]
  \quad\text{(stationary law of the closed loop).}
\]
\textbf{Role.}  Converts trajectory optimization into optimization of a
stationary distribution.  Every later formula computes stationary
expectations.\\
\textbf{Defect.}  Real networks switch regimes: a host moves from wired
to radio, a radio environment degrades, a route changes.  Statistics
estimated in one regime are wrong in the next.  The size of the defect is
the loss incurred during adaptation: (regime switch rate) $\times$
(adaptation time) $\times$ (excess loss during adaptation).  The
adaptation time of the estimators is a design parameter
(\cref{sec:controller}); the derivation does not hide it.\\
\textbf{Monitor.}  Change detection on the estimator inputs; the
telemetry flags disagreement between short-window and long-window
estimates of $(\sigma_i, \lambda_i, F_i)$.
\end{simplification}

\subsection{Mean and difference coordinates}\label{sec:coordinates}

\begin{definition}[Fleet coordinates]\label{def:coordinates}
Define the fleet mean and the deviations
\[
  \Lbar(t) = \frac1N\sum_{i=1}^N L_i(t),
  \qquad
  \Ltil{i}(t) = L_i(t) - \Lbar(t),
  \qquad
  \textstyle\sum_i \Ltil{i}(t) = 0 ,
\]
and the same for warp: $\bar u = \frac1N\sum_i u_i$,
$\tilde u_i = u_i - \bar u$.
\end{definition}

\begin{proposition}[The loss in fleet coordinates]\label{prop:loss-coords}
The pointwise loss \cref{eq:pointwise-loss} equals
\begin{equation}\label{eq:loss-coords}
  \Lambda
  \;=\;
  \underbrace{\alpha \Lbar \;+\; r\,\dot{\bar u}^2}_{\text{mean
  coordinate}}
  \;+\;
  \underbrace{\frac{q}{N}\sum_i \Ltil{i}^2
  \;+\; \frac{r}{N}\sum_i \dot{\tilde u}_i^2}_{\text{difference
  coordinates}}
  \;+\;
  \underbrace{\frac{c}{N}\sum_i \ind{m_i < 0}}_{\text{couples both}} .
\end{equation}
\end{proposition}

\begin{proof}
The latency and synchronization terms are already in coordinate form by
\cref{def:coordinates}.  For the warp term, expand
$\dot u_i = \dot{\bar u} + \dot{\tilde u}_i$:
\[
  \frac1N\sum_i \dot u_i^2
  = \frac1N\sum_i \bigl(\dot{\bar u} + \dot{\tilde u}_i\bigr)^2
  = \dot{\bar u}^2
    + \frac{2\dot{\bar u}}{N}\sum_i \dot{\tilde u}_i
    + \frac1N\sum_i \dot{\tilde u}_i^2
  = \dot{\bar u}^2 + \frac1N\sum_i \dot{\tilde u}_i^2 ,
\]
because $\sum_i \tilde u_i = 0$ identically, hence
$\sum_i \dot{\tilde u}_i = 0$.  The indicator term depends on
$m_i = \Lbar + \Ltil{i} - D_i(s_i(t))$, which contains both coordinates;
it does not decompose.
\end{proof}

\begin{proposition}[Dynamics decouple]\label{prop:dynamics-decouple}
From \cref{prop:latency-dynamics}, by averaging and subtracting:
\[
  \dot{\Lbar} = \bar w - \bar u ,
  \qquad
  \dot{\Ltil{i}} = \tilde w_i - \tilde u_i .
\]
The mean coordinate is driven only by mean drift and mean warp; each
difference coordinate only by its own drift and warp deviations.
\end{proposition}

\begin{proof}
Linearity of \cref{eq:latency-dynamics} in $(w_i, u_i)$ and of the
averaging map.
\end{proof}

\begin{remark}[The obstacle, stated precisely]\label{rem:obstacle}
After \cref{prop:loss-coords,prop:dynamics-decouple}, the problem would
split into $1 + (N{-}1)$ independent problems --- one for the mean, one
per difference direction --- \emph{except} for the indicator term, which
evaluates the sum $\Lbar + \Ltil{i}$ against the delay realization.
Everything in \cref{sec:indicator} exists to handle this one term; the
result (\cref{sec:composition}) is that after quantified approximations
it attaches only to the mean coordinate.
\end{remark}

\subsection{Superposition}\label{sec:superposition}

\begin{proposition}[Component superposition]\label{prop:superposition}
Under any linear time-invariant policy, every closed-loop signal (latency
error, margin, warp) is a sum of its responses to $B_i$, $\varepsilon_i$,
$P_i$, drift $w_i$, and measurement noise $n_i$ separately.
\end{proposition}

\begin{proof}
The closed loop of a linear plant and a linear policy is a linear map of
its inputs; \cref{as:decomposition} writes the disturbance as a sum.
\end{proof}

This licenses the accounting of \cref{sec:transfer}: each loss term is
evaluated component by component.  Cross terms between components appear
only in second moments of sums; \cref{simp:independence} (S9) treats
them.

%% =============================================================================
\section{The Quadratic Theory}\label{sec:quadratic}

This section solves the part of the problem that lives in the closed
algebra of linear dynamics, quadratic costs, and second moments.  The
solution machinery is classical; the derivation is nevertheless complete,
because the specific constants matter later.

\subsection{The closed algebra}\label{sec:closure}

Three facts, jointly, make the quadratic part solvable in closed form.

\begin{proposition}[Closure]\label{prop:closure}
\begin{enumerate}[label=(\roman*)]
  \item A linear map of a Gaussian vector is Gaussian (imported; standard
    property of Gaussian measures).
  \item The expectation of a quadratic form depends only on the first two
    moments (\cref{prop:quadratic-second-moment}).
  \item The minimizer of a quadratic cost under linear dynamics is linear
    state feedback (proven for our system in \cref{prop:riccati} below).
\end{enumerate}
Consequently: starting from linear dynamics, quadratic costs, and a
disturbance described by second moments, every object generated by the
optimization (optimal law, closed loop, cost value) stays inside a family
parameterized by matrices.  No path-space computation is needed.
\end{proposition}

\begin{remark}[Where Gaussianity actually matters]
\label{rem:gaussian-role}
For the fully observed problem \cref{eq:lqr-setup} below, the
completion-of-squares argument of \cref{prop:riccati} never uses
Gaussianity: the linear law is optimal among \emph{all} causal
state-feedback laws for any mean-zero white noise of finite intensity,
because the quadratic value function ansatz closes regardless of the
noise law.  Gaussianity earns its keep one level up: when the
controlled error must be inferred from a corrupted measurement, the
linear estimator is optimal among all causal estimators only for
Gaussian corruption.  The real assumption of the fast layer is
therefore about \emph{measurement}, and \cref{simp:gaussian} states it
as such.
\end{remark}

(Simplifications are numbered by their position in the register of
\cref{sec:register}, not by order of appearance: S10, S12, and S13
arise inside the machinery and composition chapters, S2--S9 in the
indicator chapter.)

\setcounter{simplification}{9}%
\begin{simplification}[S10: full-information law on the raw
measurement]\label{simp:gaussian}
\textbf{Statement.}  The fast layer applies the full-information
optimal law of \cref{prop:riccati} with the measured output
substituted for the true error --- no filtering stage --- and stays in
the linear time-invariant class.\\
\textbf{Role.}  Keeps the fast layer a two-gain law; avoids carrying
an estimator design through the derivation.  Under exact error
observation the restriction is vacuous (\cref{rem:gaussian-role}).\\
\textbf{Defect.}  Two parts, by backend.  (a) Under the latency
measurement, the corruption $n_i$ enters the loop and moves true
latency with transfer $-\kappa^2/\mathcal{D}(s)$
(\cref{prop:transfer}); the exposure is
$\int |\kappa^2/\mathcal{D}(j\omega)|^2 S_n(\omega)\,
\dd\omega/2\pi$, which also upper-bounds what an optimal filtering
stage could recover.  On time-disciplined hosts $n_i$ is microseconds
and the forfeit is negligible; on an undisciplined host $n_i$ drifts
by milliseconds and no linear filter repairs it --- the correct
response is switching that stream's backend, not a filter.  (b) Under
the margin measurement, the deviation itself drives the loop, and its
distribution is far from Gaussian (the pause component): a nonlinear
law can strictly beat the linear one.  The event-gated machinery of
the implementation (detection, freezing, supervision) is exactly such
a nonlinear element, added outside this derivation; its benefit is
bounded by the pause leakage terms of \cref{prop:transfer}.\\
\textbf{Monitor.}  Innovation histogram and spectrum of the loop;
measured $S_n$ level per host.
\end{simplification}

\subsection{Spectral evaluation}\label{sec:spectral}

\begin{proposition}[Output spectra of linear systems]\label{prop:psd}
Let $x$ be stationary with power spectral density $S_x(\omega)$, let $h$
be the impulse response of a stable linear time-invariant system with
transfer function $H(s)$, and let $y = h * x$.  Then $y$ is stationary
with
\begin{equation}\label{eq:psd-rule}
  S_y(\omega) = \bigl|H(j\omega)\bigr|^2 S_x(\omega),
  \qquad
  \Var[y] = \int_{-\infty}^{\infty} |H(j\omega)|^2 S_x(\omega)
  \,\frac{\dd\omega}{2\pi}.
\end{equation}
\end{proposition}

\begin{proof}
The autocovariance of $y$ is
$R_y(\tau) = \E[y(t)y(t{+}\tau)]
= \iint h(a)h(b)\,R_x(\tau + a - b)\dd a\dd b$.
Taking Fourier transforms in $\tau$ and using the convolution theorem
twice gives $S_y = H(j\omega)\overline{H(j\omega)}\,S_x$.  Integrating
the spectral density gives the variance (Bochner / Wiener--Khinchin,
imported).
\end{proof}

\begin{convention}[Transform notation]\label{conv:transforms}
In spectral and transfer-function contexts, $H(s)$ denotes a generic
transfer function and $s$ the Laplace variable; capital letters
($E, U, W, Y, M, N_i, \tilde D$) denote transforms of the
corresponding time-domain signals.  Neither use collides with stream
position $s$ (\cref{def:stream}) or episode heights $H_{i,k}$
(\cref{as:pause}), which never appear inside transforms.
\end{convention}

\subsection{The fast loop is a linear-quadratic regulator}\label{sec:lqr}

Consider one difference coordinate (\cref{sec:speaker-modes} shows all
$N{-}1$ of them are identical copies).  Write $e$ for the coordinate's
latency error, $u$ for its warp, $v = \dot u$ for the control input, and
let white noise of intensity $\sigma_w^2$ drive $e$ (per
\cref{cor:no-brownian}, this white noise represents baseline-wander and
measurement-noise components; the stationary diffusion enters separately
through \cref{prop:psd} and the transfer functions of
\cref{sec:transfer}).  The system and cost are
\begin{equation}\label{eq:lqr-setup}
  \frac{\dd}{\dd t}\begin{pmatrix} e\\ u\end{pmatrix}
  =
  \underbrace{\begin{pmatrix} 0 & -1\\ 0 & 0\end{pmatrix}}_{A}
  \begin{pmatrix} e\\ u\end{pmatrix}
  +
  \underbrace{\begin{pmatrix} 0\\ 1\end{pmatrix}}_{B} v
  + \begin{pmatrix} \sigma_w \dot W\\ 0\end{pmatrix},
  \qquad
  \text{cost rate } q e^2 + r v^2 .
\end{equation}

\setcounter{simplification}{11}%
\begin{simplification}[S12: white-noise surrogate for the loop-driving
disturbance]\label{simp:whitenoise}
\textbf{Statement.}  The disturbance driving the fast-loop design is
modeled as white noise of intensity $\sigma_w^2$, although the physical
components it stands for --- the wander of $B_i$
(\cref{as:baseline}: bounded derivative, not integrated white noise)
and measurement error --- have non-flat spectra.\\
\textbf{Role.}  Yields the closed-form Riccati and Lyapunov results
below.\\
\textbf{Defect.}  Two distinct effects.  For the \emph{gains}: the
optimal linear law against a colored disturbance augments the design
with a disturbance model; for the dominant mismatch here --- the
near-DC drift component --- that augmentation is exactly the integral
action of \cref{sec:integral}, which the implementation carries, so
the residual gain defect is confined to mid-band spectral shape.  For
the \emph{evaluations}: the variance formulas of
\cref{prop:lyapunov} err by
$\int |H(j\omega)|^2\,(S_{\mathrm{true}}(\omega) - \sigma_w^2)
\,\dd\omega/2\pi$ per output, a measurable quantity.\\
\textbf{Monitor.}  Flatness of the loop-innovation spectrum within the
bandwidth $\kappa$.
\end{simplification}

\begin{proposition}[Riccati solution]\label{prop:riccati}
The average-cost optimal law for \cref{eq:lqr-setup} is
\begin{equation}\label{eq:lqr-law}
  v \;=\; \kappa^2\, e \;-\; \sqrt2\,\kappa\, u,
  \qquad
  \kappa \;\coloneqq\; \Bigl(\frac{q}{r}\Bigr)^{1/4},
\end{equation}
with closed-loop characteristic polynomial
$\mathcal{D}(s) = s^2 + \sqrt2\,\kappa s + \kappa^2$: natural frequency
$\omega_n = \kappa$ and damping ratio $\zeta = 1/\sqrt2$.  The gains do
not depend on the noise intensity $\sigma_w^2$.
\end{proposition}

\begin{proof}
Seek a value function $V(x) = x^\top P x$ and average cost
$\varrho$ (the average-cost constant; unrelated to the recovery time
$\rho$)
satisfying the average-cost Hamilton--Jacobi--Bellman equation
\[
  \varrho \;=\; \min_v\Bigl[\, q e^2 + r v^2
    + \nabla V^\top (Ax + Bv)\Bigr]
    + \tfrac12 \sigma_w^2\,\partial^2_{ee} V .
\]
(The verification theorem for average-cost linear-quadratic control is
imported; the computation below exhibits the solution it certifies.)
The minimization over $v$ is a quadratic:
$\partial_v[\,rv^2 + 2x^\top P B v\,] = 0$ gives
$v^* = -\tfrac1r B^\top P x$.  Substituting back, the $x$-dependent part
must vanish for all $x$, giving the algebraic Riccati equation
\[
  A^\top P + PA - \tfrac1r PBB^\top P + Q = 0,
  \qquad Q = \begin{pmatrix} q&0\\0&0\end{pmatrix},
\]
and $\varrho = \sigma_w^2 P_{11}$.  Write
$P = \bigl(\begin{smallmatrix} p_{11}&p_{12}\\
p_{12}&p_{22}\end{smallmatrix}\bigr)$.  Componentwise:
\[
  A^\top P + PA
  = \begin{pmatrix} 0 & -p_{11}\\ -p_{11} & -2p_{12}\end{pmatrix},
  \qquad
  \tfrac1r PBB^\top P
  = \tfrac1r\begin{pmatrix} p_{12}^2 & p_{12}p_{22}\\
     p_{12}p_{22} & p_{22}^2\end{pmatrix}.
\]
The three scalar equations are
\[
  \text{(1,1)}:\; q = \tfrac{p_{12}^2}{r},
  \qquad
  \text{(1,2)}:\; p_{11} = -\tfrac{p_{12}p_{22}}{r},
  \qquad
  \text{(2,2)}:\; -2p_{12} = \tfrac{p_{22}^2}{r}.
\]
Equation (2,2) forces $p_{12} < 0$, so (1,1) gives
$p_{12} = -\sqrt{qr}$; then (2,2) gives
$p_{22} = \sqrt{2r\sqrt{qr}} = \sqrt2\, q^{1/4} r^{3/4}$ and (1,2) gives
$p_{11} = \sqrt2\, q^{3/4} r^{1/4} > 0$.  Positive definiteness:
$p_{11} > 0$ and
$\det P = p_{11}p_{22} - p_{12}^2 = 2qr - qr = qr > 0$.
The law is $v = -\tfrac1r(p_{12} e + p_{22} u)
= \sqrt{q/r}\,e - \sqrt2\,(q/r)^{1/4} u$, which is \cref{eq:lqr-law}
with $\kappa = (q/r)^{1/4}$.  The closed-loop matrix
$A - \tfrac1r BB^\top P
= \bigl(\begin{smallmatrix} 0&-1\\ \kappa^2& -\sqrt2\kappa
\end{smallmatrix}\bigr)$
has characteristic polynomial
$s(s+\sqrt2\kappa) + \kappa^2 = s^2 + \sqrt2\kappa s + \kappa^2$, i.e.,
$\omega_n = \kappa$, $\zeta = \sqrt2\kappa/(2\kappa) = 1/\sqrt2$,
stable.  The noise intensity appears only in $\varrho$, not in the Riccati
equation, so the gains are noise-independent.
\end{proof}

\begin{remark}[Correspondence with the deployed controller]
\label{rem:deployed}
The implementation runs exactly this second-order law (plus integral
action, \cref{sec:integral}), derived earlier in sample units with cost
$\Var[e_{\mathrm{samples}}] + \lambda \Var[\dot u]$ (that
document's smoothness weight $\lambda$, unrelated to the episode
rates $\lambda_i$).  Converting
$e_{\mathrm{samples}} = F_s\, e_{\mathrm{seconds}}$ maps that cost to
$q = F_s^2$, $r = \lambda$, hence
$\kappa = (F_s^2/\lambda)^{1/4}$ --- the same natural frequency, and the
same universal damping $\zeta = 1/\sqrt2$.  What is new here is not the
law but its \emph{scope}: \cref{prop:riccati} makes it the
full-information optimum over all causal laws for the difference
coordinates --- not a choice within a two-parameter family --- and S10
states exactly what running it on the raw measurement gives up.
\end{remark}

\begin{proposition}[Stationary variances of the fast loop]
\label{prop:lyapunov}
Under the law \cref{eq:lqr-law} driven by white noise of intensity
$\sigma_w^2$ on $e$, the stationary covariance
$X = \Cov(e,u)$ solves the Lyapunov equation
$A_c X + X A_c^\top + \operatorname{diag}(\sigma_w^2, 0) = 0$ with
$A_c = \bigl(\begin{smallmatrix} 0&-1\\ \kappa^2&-\sqrt2\kappa
\end{smallmatrix}\bigr)$, and
\begin{equation}\label{eq:lyapunov-values}
  \Var[e] = \frac{3\,\sigma_w^2}{2\sqrt2\,\kappa},
  \qquad
  \Cov[e,u] = \frac{\sigma_w^2}{2},
  \qquad
  \Var[u] = \frac{\kappa\,\sigma_w^2}{2\sqrt2},
  \qquad
  \Var[\dot u] = \frac{\kappa^3 \sigma_w^2}{2\sqrt2}.
\end{equation}
\end{proposition}

\begin{proof}
Write $X = \bigl(\begin{smallmatrix} x_e&x_c\\ x_c&x_u
\end{smallmatrix}\bigr)$.  Compute
$A_c X = \bigl(\begin{smallmatrix} -x_c & -x_u\\
\kappa^2 x_e - \sqrt2\kappa x_c & \kappa^2 x_c - \sqrt2\kappa x_u
\end{smallmatrix}\bigr)$.
The Lyapunov equation componentwise:
\[
  \text{(1,1)}:\; -2x_c + \sigma_w^2 = 0
  \;\Rightarrow\; x_c = \tfrac{\sigma_w^2}{2};
  \qquad
  \text{(2,2)}:\; 2\kappa^2 x_c - 2\sqrt2\kappa x_u = 0
  \;\Rightarrow\; x_u = \tfrac{\kappa x_c}{\sqrt2}
   = \tfrac{\kappa\sigma_w^2}{2\sqrt2};
\]
\[
  \text{(1,2)}:\; -x_u + \kappa^2 x_e - \sqrt2\kappa x_c = 0
  \;\Rightarrow\;
  x_e = \frac{x_u + \sqrt2\kappa x_c}{\kappa^2}
      = \frac{\sigma_w^2}{\kappa}
        \Bigl(\frac{1}{2\sqrt2} + \frac{\sqrt2}{2}\Bigr)
      = \frac{3\,\sigma_w^2}{2\sqrt2\,\kappa}.
\]
For $\Var[\dot u]$, use $\dot u = v = \kappa^2 e - \sqrt2 \kappa u$:
\[
  \Var[v] = \kappa^4 x_e - 2\sqrt2\kappa^3 x_c + 2\kappa^2 x_u
  = \kappa^3\sigma_w^2\Bigl(\frac{3}{2\sqrt2} - \sqrt2
    + \frac{1}{\sqrt2}\Bigr)
  = \frac{\kappa^3\sigma_w^2}{2\sqrt2},
\]
using $\tfrac{3}{2\sqrt2} - \sqrt2 + \tfrac1{\sqrt2}
= \tfrac{3 - 4 + 2}{2\sqrt2} = \tfrac1{2\sqrt2}$.
\end{proof}

\subsection{Speaker-space diagonalization}\label{sec:speaker-modes}

\begin{proposition}[Mode decomposition]\label{prop:modes}
Let $\Pi = I_N - \tfrac1N \mathbf{1}\mathbf{1}^\top$ be the centering
matrix.  Then:
\begin{enumerate}[label=(\roman*)]
  \item The synchronization cost is
    $\tfrac{q}{N} L^\top \Pi L$; the matrix $\Pi$ has eigenvalue $0$ on
    the all-ones direction (the mean coordinate) and eigenvalue $1$ on
    its orthogonal complement (the difference coordinates).
  \item In any orthonormal basis whose first vector is
    $\mathbf{1}/\sqrt N$, the quadratic part of the loss and the
    dynamics \cref{prop:dynamics-decouple} split into: a mean-coordinate
    problem with \emph{zero} quadratic state cost, linear cost
    $\alpha\Lbar$, and warp-rate cost; and $N-1$ identical difference
    problems, each an instance of \cref{eq:lqr-setup} with weights
    $(q/N, r/N)$ --- the common factor $1/N$ cancels in the ratio
    $q/r$, so the bandwidth $\kappa$ of \cref{prop:riccati} is
    unaffected and the factor is suppressed from here on.
\end{enumerate}
\end{proposition}

\begin{proof}
(i) $\Pi\mathbf{1} = \mathbf{1} - \mathbf{1} = 0$; for $x \perp
\mathbf{1}$, $\Pi x = x$.  Also
$\sum_i (L_i - \Lbar)^2 = L^\top \Pi L$ by expanding the square.
(ii) The warp cost $\tfrac{r}{N}\sum \dot u_i^2 = \tfrac{r}{N}
\dot u^\top I \dot u$ is proportional to the identity, hence diagonal in
every orthonormal basis; the dynamics are the same scalar integrator in
every coordinate (\cref{prop:dynamics-decouple}); the synchronization
cost is diagonal by (i) with the stated eigenvalues; the latency cost
$\alpha\Lbar$ involves only the first basis vector.  Disturbance
covariances need not be diagonal in this basis; by
\cref{prop:riccati} the optimal gains do not depend on them, and cost
evaluation projects the covariances onto the modes.
\end{proof}

\begin{remark}[What the mean coordinate is left with]
\label{rem:mean-coordinate}
The mean coordinate has no quadratic restoring cost: synchronization
puts eigenvalue zero on it, and the warp-rate cost charges only
\emph{motion} of the mean warp, not its value.  Its only costs are the
linear latency price $\alpha\Lbar$ and (through the margin) the
indicator term.  The quadratic machinery therefore has nothing to say
about \emph{where} the fleet sits --- only about how it moves.  This
degeneracy is not an accident; it is the structural hook on which
\cref{sec:composition} hangs the underrun term.
\end{remark}

\subsection{Closed-loop transfer functions and the component ledger}
\label{sec:transfer}

The fast law \cref{eq:lqr-law} in transfer form: from $s U = \kappa^2 Y
- \sqrt2\kappa U$,
\begin{equation}\label{eq:controller-tf}
  U(s) = C(s)\, Y(s),
  \qquad
  C(s) = \frac{\kappa^2}{s + \sqrt2\kappa},
\end{equation}
where $Y$ is the measured error.  The plant is $sE = W - U$ with $W$ the
drift disturbance.  Let $\tilde D$ denote the transform of the
deviation of the delay process from its baseline --- the same physical
quantity written $x_i$ in the time domain
(\cref{sec:indicator-exact}) --- $N_i$ the transform of the
measurement error $n_i$ of \cref{def:meas-e2e}, and
$\mathcal{D}(s) = s^2 + \sqrt2\kappa s + \kappa^2$.

\begin{proposition}[Backend transfer functions]\label{prop:transfer}
Under the margin measurement ($Y = E - \tilde D$, \cref{def:meas-niq}):
\begin{equation}\label{eq:tf-niq}
  \frac{E}{\tilde D} = \frac{\kappa^2}{\mathcal{D}(s)},
  \qquad
  \frac{M}{\tilde D} = -\,\frac{s\,(s+\sqrt2\kappa)}{\mathcal{D}(s)},
  \qquad
  \frac{U}{\tilde D} = -\,\frac{\kappa^2 s}{\mathcal{D}(s)},
\end{equation}
where $M = E - \tilde D$ is the margin deviation.  Under the latency
measurement ($Y = E + N_i$, \cref{def:meas-e2e}):
\begin{equation}\label{eq:tf-e2e}
  \frac{E}{N_i} = -\,\frac{\kappa^2}{\mathcal{D}(s)},
  \qquad
  \frac{E}{\tilde D} = 0,
  \qquad
  \frac{M}{\tilde D} = -1 \;\;\text{(all frequencies)} .
\end{equation}
\end{proposition}

\begin{proof}
Margin measurement: $sE = W - C\,(E - \tilde D)$ gives
$E\,(s + C) = W + C\tilde D$, so
$E/\tilde D = C/(s+C)$.  Compute
$s + C = \bigl(s(s+\sqrt2\kappa) + \kappa^2\bigr)/(s+\sqrt2\kappa)
= \mathcal{D}(s)/(s+\sqrt2\kappa)$, hence
$C/(s+C) = \kappa^2/\mathcal{D}(s)$.  Then
$M/\tilde D = E/\tilde D - 1 = (\kappa^2 - \mathcal{D}(s))/\mathcal{D}(s)
= -s(s+\sqrt2\kappa)/\mathcal{D}(s)$, and
$U = C\,(E - \tilde D) = C\,(M/\tilde D)\,\tilde D
= -\kappa^2 s/\mathcal{D}(s)\,\tilde D$.
Latency measurement: $sE = W - C(E + N_i)$ gives
$E/N_i = -C/(s+C) = -\kappa^2/\mathcal{D}(s)$; the delay deviation does not enter
the loop equations at all, so $E/\tilde D = 0$ and
$M = E - \tilde D$ carries $\tilde D$ with gain $-1$ at every
frequency.
\end{proof}

\begin{remark}[Reading \cref{prop:transfer}]\label{rem:transfer-reading}
Under the margin measurement, $E/\tilde D$ is a unit-gain low-pass:
slow delay changes (baseline) are converted one-for-one into latency
changes, and only fluctuations above the loop bandwidth $\kappa$ land in
the margin.  The pause component is far above $\kappa$ in its onset, so
a pause drives the margin almost fully --- but the loop still responds
through $U/\tilde D$: pauses leak into warp (and hence into the
synchronization and roughness terms) under this backend.
Under the latency measurement the network deviation never drives the
loop: the entire deviation, at \emph{all} frequencies, lands in the
margin.  Slow baseline growth then consumes margin silently --- the
level layer must track the baseline explicitly
(\cref{sec:level-program}) --- and a drifting clock error $n_i$ moves
true latency one-for-one at DC.  Neither backend removes a disturbance;
each chooses which loss term absorbs it.
\end{remark}

\begin{table}[h]
\centering
\caption{Which loss term absorbs which disturbance component.}
\label{tab:ledger}
\begin{tabular}{@{}lll@{}}
\toprule
Component & Latency measurement (e2e) & Margin measurement (niq)\\
\midrule
Baseline $B_i$ & indicator (silent margin drain; & latency term
  (latency follows \\
 & level layer must track $\hat B_i$) & the network) \\
Diffusion $\varepsilon_i$ & indicator (bulk tail term) & margin
  (high-pass part) $\to$ indicator; \\
 & & warp/sync (band-pass leakage) \\
Pauses $P_i$ & indicator only & indicator $+$ warp/sync leakage \\
Drift $w_i$ & warp DC (free) & warp DC (free) \\
Clock error $n_i$ & latency term (unit DC gain) & absent \\
\bottomrule
\end{tabular}
\end{table}

\subsection{Integral action}\label{sec:integral}

The drift $w_i$ is unknown and slowly varying.  The principled
treatment augments the state with the unknown input ($\dot w = 0$) and
solves the augmented problem; by the internal-model principle
(imported: Francis and Wonham, 1976), any law achieving zero
steady-state error against a constant disturbance must contain an
integrator, and the certainty-equivalent reduction of the augmented
design --- estimate $\hat w$ by integrating the loop innovation, feed
it forward --- lands on the law \cref{eq:lqr-law} plus an integral
term:
\begin{equation}\label{eq:pid}
  \dot\xi = e,
  \qquad
  v = \kappa^2 e + K_0\,\xi - \sqrt2\kappa\,u ,
\end{equation}
with $K_0$ placing a pole well below $\kappa$ (dominant-pole design).
Two caveats, so the status of \cref{eq:pid} is exact:
it is certainty-equivalent and dominant-pole-tuned, not the jointly
optimized augmented-state law ($K_0$ is set by a chosen settling time
$T_I$, not by the augmented Riccati equation); and the integral mode
inflates the stationary error variance by the factor
$1/\bigl(1 - 1/(\sqrt2\,\kappa\,T_I)\bigr) \approx 1.035$ --- a
$3.5\%$ overhead at the deployed
$\kappa \approx 0.69\,\mathrm{s}^{-1}$, $T_I = 30$\,s.  The
deployed controller implements exactly \cref{eq:pid} with anti-windup
on the warp clamp; the earlier document\footnote{\texttt{controller.tex}
in this directory: the second-order controller derivation.  Its gain
optimization is consistent with \cref{prop:riccati}; its noise passage
is superseded by \cref{prop:telescope,cor:no-brownian}: a stationary
jitter level cannot supply the accumulating white noise assumed there,
so its $\sigma_w$ must be read as baseline wander plus measurement
noise, and its Gaussian-tail underrun estimates are replaced by
\cref{sec:indicator}.}
gives the dominant-pole design, the variance-overhead computation, and
stability margins (the integral gain must satisfy
$K_0 < \sqrt2\,\kappa\cdot\kappa^2$, the Routh bound).

%% =============================================================================
\section{The Indicator Term}\label{sec:indicator}

The underrun term is outside the quadratic algebra
(\cref{sec:indicator-not-moments}).  This section evaluates it by a chain
of quantified simplifications.  The end product is a function
$R_i(\mu)$ --- underrun fraction versus margin level --- consumed by the
level program of \cref{sec:composition}.

\subsection{Exact rewriting}\label{sec:indicator-exact}

Fix receiver $i$.  Under S1, the underrun contribution to the loss is
$c$ times the stationary probability $\Prob(m_i < 0)$.  Decompose the
margin using \cref{as:decomposition}: with the fast loop holding the
latency near a slowly varying level $\ell_i(t)$,
\begin{equation}\label{eq:margin-decomp}
  m_i \;=\; \ell_i - B_i - x_i + e_i^{\mathrm{loop}},
  \qquad
  x_i \coloneqq \varepsilon_i + P_i ,
\end{equation}
where $x_i$ is the \emph{deviation} (the part of the delay above
baseline; this is the quantity the telemetry's deviation meter measures)
and $e_i^{\mathrm{loop}}$ is the fast loop's tracking error.  Define the
\emph{margin level}
\begin{equation}\label{eq:margin-level}
  \mu_i \;\coloneqq\; \ell_i - \hat B_i
  \quad\text{(latency backend)},
\end{equation}
the distance between the held latency and the estimated baseline.
Under the margin backend the loop regulates the margin directly and
$\mu_i$ is the configured queue target itself, with no baseline
subtraction (\cref{rem:backend-deviation}); every use of $\mu_i$
below is backend-conditional in this sense.  Then
\begin{equation}\label{eq:R-def}
  \Prob(m_i < 0)
  \;=\;
  \Prob\bigl(x_i > \mu_i\bigr) \;+\; \text{(errors)},
  \qquad
  R_i(\mu) \coloneqq \Prob(x_i > \mu),
\end{equation}
where the errors come from baseline estimation error
$B_i - \hat B_i$, bounded by the estimator lag $\bar\beta\,\tau_{\hat B}$
(with $\tau_{\hat B}$ the baseline filter's time constant, distinct
from the emission time $\tau(s)$), and from
$e_i^{\mathrm{loop}}$, whose stationary scale is
$\sqrt{\Var[e]} = \sqrt{3\sigma_w^2/(2\sqrt2\,\kappa)}$ by
\cref{prop:lyapunov} --- both far smaller than $\sigma_i$ at the
deployed constants.  $R_i$ is the exceedance function of
the stationary deviation law --- a directly measurable object: the
deviation histogram estimates it without any model.

\begin{remark}[Backend dependence of the deviation law]
\label{rem:backend-deviation}
\Cref{eq:margin-decomp} reads differently per backend.  Under the
latency measurement, the loop holds $L_i$ at $\ell_i$ and the full
deviation lands in the margin at every frequency
(\cref{prop:transfer}): $x_i$ enters unfiltered.  Under the margin
measurement, the loop absorbs the part of $x_i$ below the bandwidth
$\kappa$ into latency, so the deviation reaching the margin is
high-passed, and the margin level $\mu_i$ is the configured queue
target itself rather than $\ell_i - \hat B_i$ (the baseline is
absorbed by the loop and drops out of \cref{eq:margin-level}).  In
both cases the deviation histogram is measured on the running closed
loop, so $\hat R_i$ estimates the correct backend-specific exceedance
law by construction; only the component-model \emph{extrapolation} of
\cref{sec:risk-evaluation} must respect the backend.
\end{remark}

\begin{definition}[Occupation density]\label{def:occupation}
$g_i(\mu) \coloneqq -\,\dd R_i/\dd\mu$: the probability density of time
the deviation spends at level $\mu$.  The level program's optimality
condition (\cref{sec:level-program}) is stated in terms of $g_i$.
\end{definition}

\subsection{Evaluation under the component model}\label{sec:risk-evaluation}

The component model of \cref{sec:process} evaluates $R_i$ as three
terms.  The decomposition serves two purposes: it \emph{extrapolates}
$R_i$ beyond the directly observed range (the tail of the pause height
law extends past any finite observation window), and it \emph{explains}
which mechanism dominates at which margin.

\begin{proposition}[Risk evaluation]\label{prop:risk-evaluation}
Under S1--S5 below,
\begin{equation}\label{eq:risk-evaluation}
  R_i(\mu)
  \;=\;
  \underbrace{\Prob(\varepsilon_i > \mu)}_{\text{bulk}}
  \;+\;
  \underbrace{\lambda_i\,\E\pos{H_i - \mu}}_{\text{pauses
  (\cref{prop:campbell})}}
  \;+\;
  \underbrace{\lambda_i \int \E\Bigl[\pos{H' - \mu + \delta(g)}
      - \pos{H' - \mu}\Bigr]\, G_i(\dd g)}_{\text{compounding}}
  \;+\; \text{defects of S2--S5},
\end{equation}
where $G_i$ is the gap law between consecutive episodes, $H'$ is the
height of the following episode, and
\begin{equation}\label{eq:depression}
  \delta(g) \;=\; \chi\,\pos{1 - g/\rho}
\end{equation}
is the residual margin depression at gap $g$ after an episode that
consumed margin $\chi = \min(H, \mu)$ and recovers over time $\rho$
(\cref{as:recovery}; linear recovery model).  The expectation in the
compounding term is over three quantities: the predecessor's height
$H$ and recovery $\rho$ (entering through $\delta$), and the
follower's height $H'$; per S5 they are taken mutually independent
and independent of the gap $g$, each with its marginal law.
\end{proposition}

The bulk term is the stationary exceedance of the diffusion --- exact
under S1 outside episodes.  \textbf{With wired-host $\sigma_i$ up to one
millisecond and margin levels of order $5$--$10$\,ms, the bulk term is
comparable to or larger than the pause term}; no negligibility is
claimed anywhere.  The pause term is \cref{prop:campbell}.  The
compounding term prices an episode that arrives while the margin is
still depressed: the follower's effective margin is
$\mu - \delta(g)$ instead of $\mu$.  With fast redelivery
($\rho \sim$ milliseconds) the integral is negligible because
$G_i(\rho) \approx 0$; with warp-limited recovery
($\rho$ up to $\chi/u_{\max}$, i.e., seconds) and the measured gap
clustering (short gaps two to three orders above the exponential
prediction), the term is first-order.  It is computable from measured
objects: the gap histogram, the height law, and the recovery mechanism.

\begin{definition}[Belief state]\label{def:belief}
The \emph{belief state} $\hat\theta(t)$ is the vector of estimator
outputs feeding the risk evaluation and, later, the level program:
\[
  \hat\theta \;\coloneqq\;
  \bigl(\hat B_i,\ \hat\sigma_i,\
  \text{deviation histogram},\
  \hat\lambda_i,\ \hat F_i,\ \hat G_i,\ \hat\rho_i
  \bigr)_{i=1,\dots,N}.
\]
Wherever a later statement says ``the belief moves,'' it means these
estimates change.
\end{definition}

\subsection{The simplification chain}\label{sec:simplification-chain}

\setcounter{simplification}{1}%
\begin{simplification}[S2: additive component evaluation]
\label{simp:additive}
\textbf{Statement.}  $R_i$ is evaluated as the sum
\cref{eq:risk-evaluation}: bulk and pause exceedances add, and the
diffusion is ignored \emph{inside} episodes.\\
\textbf{Role.}  Lets each term be estimated by its own estimator (the
histogram for the bulk; the event meter for pauses).\\
\textbf{Defect.}  Three parts.  (a) \emph{Threshold smearing}: inside an
episode the true condition is $\varepsilon_i + P_i > \mu$, not
$P_i > \mu$; because the ramp crosses level $\mu$ at unit rate, the
diffusion shifts the crossing time by $O(|\varepsilon_i|)$, so the error
per episode is at most $\E|\varepsilon_i|$ seconds of misclassified
time, and the rate error is at most
$\lambda_i\,\bar F_i(\mu - O(\sigma_i))\cdot O(\sigma_i)$ relative to a
pause term whose scale is $\lambda_i \E\pos{H_i-\mu}$: the relative
defect is $O(\sigma_i / \text{height spread})$, a few percent when
heights are tens of milliseconds, tens of percent when the height law
concentrates near the margin --- in which case the bulk and pause terms
have merged physically and the histogram estimate of $R_i$ (which never
used the decomposition) remains the ground truth.  (b) \emph{Duty
cycle}: the bulk term uses all time rather than time outside episodes;
the error is a factor $1 - \lambda_i\E[H_i]$ --- order $10^{-3}$ on
wired hosts, up to order $10^{-2}$ on radio hosts with second-scale
heights.  (c) \emph{Zero-drain ramp} (\cref{prop:ramp}): per-episode
timing error bounded by the burst drain time of
\cref{as:recovery}(a).\\
\textbf{Monitor.}  Compare the composed $R_i$ of
\cref{eq:risk-evaluation} against the raw deviation-histogram exceedance
on the overlap range; disagreement is exactly the S2 defect.
\end{simplification}

\begin{simplification}[S3: the fast layer cannot move the indicator]
\label{simp:actuation}
\textbf{Statement.}  The underrun term is treated as independent of the
fast-loop policy: only the level path $\ell_i(t)$ enters it.\\
\textbf{Role.}  This is the load-bearing decoupling.  It removes the
indicator from the fast layer's optimization, which is what allows
\cref{prop:riccati} to stand as the fast-layer optimum.  Without it,
the barrier term couples to the loop and the problem is a free-boundary
dynamic program.\\
\textbf{Defect.}  During an episode, the most any admissible policy can
add to the margin before it empties is
$u_{\max} \times (\text{time to empty}) = u_{\max}\,\mu$ (playing
slower stretches $\mu$ buffered seconds into $\mu/(1-u_{\max})$ wall
seconds).  With $u_{\max} = 5\times10^{-3}$ and $\mu = 10$\,ms this is
$50\,\mu$s; with $\mu = 100$\,ms, $0.5$\,ms.  The loss difference
between the best fast policy and a fast policy that ignores the
indicator is therefore at most
\[
  \frac{c}{N}\sum_i\bigl[R_i(\mu_i) - R_i\bigl(\mu_i(1{+}u_{\max})\bigr)\bigr]
  \;\le\;
  \frac{c}{N}\sum_i g_i(\mu_i)\; u_{\max}\,\mu_i ,
\]
a relative defect of $u_{\max}\,\mu_i\,g_i(\mu_i)/R_i(\mu_i)$.  The
honest size of this ratio differs by component.  For the pause part,
$g/R$ is of order one over the height spread, so the ratio is
$u_{\max}\mu/(\text{height spread})$ --- well below one percent at the
reference numbers.  For the bulk part, $g/R$ is the tail hazard rate
of $\varepsilon_i$, of order $1/x_0$ with $x_0 \le \sigma_i$ the tail
scale: the ratio is $u_{\max}\mu/x_0$, which at $\mu = 10$\,ms and
$x_0 = 1$\,ms is $5\%$, and larger for faster-decaying tails.  The
decoupling nevertheless stands, on an absolute anchor: at the
operating point the first-order condition (\cref{prop:foc}) pins
$\tfrac{c}{N}\sum_i g_i(\mu_i) = \alpha$, so the absolute loss the
fast layer could recover is at most
$c\,\bar g\,u_{\max}\mu = \alpha\,u_{\max}\mu$ --- a fraction
$u_{\max}\,\mu/\ell \le u_{\max} = 0.5\%$ of the latency term itself.
No admissible fast policy moves the indicator term by more than
that.
The same authority argument covers reacting to bulk fluctuations: a
fluctuation of duration $T_f \sim 10$\,ms grants
$u_{\max}T_f = 50\,\mu\mathrm{s} \ll \sigma_i$.\\
\textbf{Monitor.}  Both forms: the relative ratio
$u_{\max}\mu\,g(\mu)/R(\mu)$ per component, and the absolute anchor
$\alpha\,u_{\max}\mu$ against the running loss estimate; if warp
authority is ever raised substantially, these monitors flag the
breakdown of the decomposition (\cref{sec:limits}).
\end{simplification}

\begin{simplification}[S4: frozen margin during an episode]
\label{simp:frozen}
\textbf{Statement.}  The evaluation \cref{eq:campbell} holds the margin
constant during an episode: no mid-episode warp response and no
mid-episode refill.\\
\textbf{Role.}  Gives the closed form $\pos{H-\mu}$ for per-episode
underrun time.\\
\textbf{Defect.}  Mid-episode warp is bounded by S3
($\le u_{\max}\mu$); mid-episode arrivals are excluded by the
definition of an episode (delivery stopped).  The defect is the S3
number again; it is not counted twice in the register.\\
\textbf{Monitor.}  Same as S3.
\end{simplification}

\begin{simplification}[S5: one-step compounding]
\label{simp:compound}
\textbf{Statement.}  The compounding term in \cref{eq:risk-evaluation}
considers only consecutive episode pairs, with the follower's height
drawn independently of the gap; chains of three or more episodes inside
one recovery window are dropped, as is gap--height dependence.\\
\textbf{Role.}  Keeps the compound term a one-dimensional integral over
the measured gap histogram.\\
\textbf{Defect.}  Chain terms are
$O\bigl((\lambda_i^{\mathrm{loc}}\rho)^2\bigr)$ where
$\lambda_i^{\mathrm{loc}}$ is the \emph{local} episode rate inside a
cluster (the measured short-gap rate, not the global average) ---
with strong clustering this is the first quantity to watch; gap--height
dependence is measurable from the joint (gap, height) records and enters
as a covariance correction.\\
\textbf{Monitor.}  Fraction of episodes whose predecessor lies within
$\rho$; the drain validator (windowed margin minimum), which sees
compound drains directly and arbitrates the whole of
\cref{eq:risk-evaluation} against reality.
\end{simplification}

\begin{simplification}[S6: two-timescale separation]
\label{simp:twotimescale}
\textbf{Statement.}  The level $\ell_i(t)$ and the belief
$\hat\theta(t)$ move slowly relative to episodes and to the fast-loop
bandwidth: the evaluation treats $\mu_i$ as constant within episodes
and within estimation windows.\\
\textbf{Role.}  Turns the indicator into a \emph{static} function
$R_i(\mu_i)$ of the current level; the level path enters only through
the glide cost (\cref{sec:glide}).\\
\textbf{Defect.}  During a level glide at rate $\gamma$, the margin
during an episode of duration $H$ differs from the static value by up
to $\gamma H$; relative defect $\gamma H / \sigma$-scale.  With the
glide bounded by $u_{\max}$ and $H \le 1$\,s this is at most
$5$\,ms --- not negligible during fast raises; the raise policy of
\cref{sec:glide} therefore counts the transient explicitly rather than
relying on S6.\\
\textbf{Monitor.}  Glide rate telemetry; episodes stamped with the
concurrent glide rate.
\end{simplification}

\begin{simplification}[S7: certainty equivalence]
\label{simp:certainty}
\textbf{Statement.}  The level program optimizes against the
\emph{estimated} risk functions $\hat R_i$ as if they were true; no
hedging against estimation error and no active probing.\\
\textbf{Role.}  Replaces a Bayesian dynamic program over beliefs with a
plug-in static program.\\
\textbf{Defect.}  Second-order in the estimation error: for the scalar
level program with optimum $\mu^*$, the excess loss is
$\tfrac12 J''(\mu^*)\,\E[(\hat\mu^* - \mu^*)^2]$ with curvature
$J''(\mu^*) = -\tfrac{c}{N}\sum_i g_i'(\mu_i^*)$, positive under S8's
decreasing densities; the estimator variance follows from event-count
statistics: with $n$ exceedance observations near the operating point,
$\operatorname{sd}(\hat\mu^*) \propto 1/\sqrt n$ with a clustering
inflation factor read from the gap law (\cref{rem:division}).  Active
probing is unnecessary because every episode's height is measured
whether or not it causes underrun: the risk curve at levels above the
operating point is identified from the same records, without ever
running the system at those levels.\\
\textbf{Monitor.}  Confidence band of $\hat R_i$ propagated to a
confidence interval on $\mu^*$; the band width is the S7 budget.
\end{simplification}

\begin{simplification}[S8: smooth, convex risk curve]
\label{simp:smooth}
\textbf{Statement.}  $R_i$ is differentiable at the operating point
($g_i$ exists), and the composed level objective $J$ is convex on the
operating range.\\
\textbf{Role.}  Lets the level program use the first-order condition
and report two-sided sensitivities.\\
\textbf{Defect.}  Three ways the statement can fail, none of which
breaks the program --- it is one-dimensional and solved by direct
scan of the measured $\hat J$ when the assumption fails
(\cref{prop:risk-shape}(ii)).  (a) An atom in the height law (a
mechanism producing a fixed-height pause) makes $R_i$ kinked: the
first-order condition becomes a subgradient inclusion and the optimum
may sit exactly at the atom; sensitivities become one-sided.
(b) Convexity of the bulk term requires a decreasing density beyond
the operating point (true for unimodal light-tailed laws); the pause
term is convex unconditionally
($\dd^2/\dd\mu^2\,\lambda\E\pos{H-\mu} = \lambda f_H(\mu) \ge 0$).
(c) The compounding term carries a concave piece proportional to the
residual-depression fraction (\cref{prop:risk-shape}(i) and its
proof): in the clustered, slow-recovery regime the composed $J$ can
lose convexity, and the first-order characterization with it.\\
\textbf{Monitor.}  The estimated $\hat g_i$ near the operating point;
mass concentration at single heights; the sign of the second
difference of the scanned $\hat J$ across the operating range.
\end{simplification}

\begin{simplification}[S9: component independence in second-order
accounting]\label{simp:independence}
\textbf{Statement.}  Cross-covariances between $B_i$, $\varepsilon_i$,
$P_i$ (and across receivers, between idiosyncratic parts) are dropped
where second moments are summed.\\
\textbf{Role.}  Keeps the spectral accounting of \cref{sec:transfer}
per-component.\\
\textbf{Defect.}  Measured violations exist: episodes co-move with
baseline shifts (a radio environment change moves both).  The dropped
terms are covariances bounded by the product of component scales; the
affected outputs are the quadratic terms' \emph{evaluations} (not the
gains, which are covariance-independent by \cref{prop:riccati}, and
not the risk curve, which is estimated on the summed deviation
directly).\\
\textbf{Monitor.}  Cross-correlation of the baseline estimator's drift
with event occurrence.
\end{simplification}

\subsection{Shape of the risk curve}\label{sec:risk-shape}

\begin{proposition}[Monotonicity, convexity, interiority]
\label{prop:risk-shape}
$R_i$ is nonincreasing.  Consider the composed level objective over
the fleet level $\ell$,
$J(\ell) = \alpha\ell + \tfrac{c}{N}\sum_i R_i(\ell - b_i)$ (with
fixed baseline offsets $b_i$), on the domain
$\ell \ge \ell_{\mathrm{floor}} \coloneqq \max_i b_i$ (below the
floor, some receiver's margin is negative --- permanent underrun ---
and such levels are excluded from the program).  Then:
\begin{enumerate}[label=(\roman*)]
  \item The bulk and pause contributions to each $R_i$ are convex
    under S8's density conditions; the compounding contribution is
    \emph{not convex in general} (see the proof), so convexity of the
    composed $J$ on the operating range is an \emph{assumption},
    carried by S8 and monitored, not a theorem.
  \item $J$ is continuous and $J(\ell) \ge \alpha\ell \to \infty$, so
    a minimizer exists on $[\ell_{\mathrm{floor}}, \infty)$ regardless
    of convexity; because the program is one-dimensional over a
    measured function, the minimizer is computable by direct scan of
    $\hat J$, and \cref{eq:foc} characterizes it only where $J$ is
    smooth and locally convex.
  \item Under S8 (including composed convexity), the minimizer is
    unique and is interior --- so that the first-order condition of
    \cref{prop:foc} applies --- exactly when
    \begin{equation}\label{eq:interior-condition}
      \alpha \;<\; \frac{c}{N}\sum_{i=1}^N
        g_i\bigl(\ell_{\mathrm{floor}} - b_i\bigr)\,,
    \end{equation}
    the risk relief available at the floor (each receiver evaluated
    at its own floor margin $\ell_{\mathrm{floor}} - b_i \ge 0$);
    otherwise the optimum sits at the floor --- a corner worth an
    alarm, not a controller.
\end{enumerate}
\end{proposition}

\begin{proof}
Monotonicity: exceedance functions are nonincreasing.  Convexity of
the bulk part where the density is decreasing; of the pause part
because $\dd^2/\dd\mu^2\,\E\pos{H-\mu} = f_H(\mu) \ge 0$ by
differentiating $\E\pos{H-\mu} = \int_\mu^\infty \bar F(h)\dd h$
twice.  The compounding part fails the same argument: its threshold
for the follower is $t(\mu) = \mu - c_g\min(H,\mu)$ with
$c_g = \pos{1 - g/\rho}$ the residual-depression fraction, so
$t'(\mu)$ jumps \emph{up} by $c_g$ at $\mu = H$ (from $1-c_g$ to
$1$), the slope $-t'(\mu)\bar F_{H'}(t(\mu))$ jumps down, and mixing
over the predecessor height contributes the negative (concave) piece
$-c_g\,\bar F_{H'}\bigl((1-c_g)\mu\bigr) f_H(\mu)$ to $J''$ alongside
the positive $t'^2 f_{H'}$ pieces.  The concave piece scales with
$c_g$: it is largest exactly in the clustered, slow-recovery regime.
Hence (i).  Existence by continuity and coercivity upward; the scan
statement is immediate from one-dimensionality.  Interiority:
$J'(\ell_{\mathrm{floor}}^+) = \alpha -
\tfrac{c}{N}\sum_i g_i(\ell_{\mathrm{floor}} - b_i)$; under composed
convexity, \cref{eq:interior-condition} makes this negative and the
minimizer interior; if it fails, $J' \ge 0$ on the whole domain and
the floor is optimal.
\end{proof}

%% =============================================================================
\section{Composition: the Two Solvers and the Free Direction}
\label{sec:composition}

The quadratic machinery (\cref{sec:quadratic}) and the risk evaluation
(\cref{sec:indicator}) use disjoint mathematics and disjoint statistics.
This section shows why they compose without conflict, derives the level
program, and states the decomposition theorem.

\subsection{The free direction}\label{sec:free-direction}

Collect three facts already proven:
\begin{enumerate}[label=(F\arabic*)]
  \item The synchronization cost has eigenvalue zero on the mean
    coordinate (\cref{prop:modes}): moving all receivers together is
    free for synchronization.
  \item The warp cost charges only the \emph{rate of change} of warp:
    a held latency level costs nothing, and a steady glide of the level
    costs nothing either --- during a constant-rate glide the
    commanded mean warp is constant, so its contribution to
    $\dot{\bar u}$ is zero.  Only changes of the glide rate are
    charged, through $\dot{\bar u} = -\ddot\ell$ up to two residuals
    --- drift-rate variation and the common-mode part of the fast
    loops' disturbance response --- both priced by S13 below
    (\cref{prop:loss-coords,prop:dynamics-decouple}).
  \item After S2--S6, the indicator term depends on the policy only
    through the slowly varying margin levels
    $\mu_i(t) = \ell_i(t) - \hat B_i(t)$
    (\cref{sec:risk-evaluation}).
\end{enumerate}
The quadratic part of the loss is thus \emph{degenerate} exactly on the
coordinate --- the static mean level --- through which the indicator
term acts.  The two solvers meet on orthogonal coordinates: the
quadratic solver fills every mode and every frequency except the mean
coordinate's zero-frequency component; the risk program occupies
precisely that component.  The loss was shaped for this: the linear
latency term (no quadratic pull toward any preset level), the
variance-only synchronization term (no absolute reference), and the
derivative-only warp cost (no charge for holding an offset) each keep
the static mean level free of quadratic forces.

\setcounter{simplification}{12}%
\begin{simplification}[S13: common-mode loop response priced
separately]\label{simp:commonmode}
\textbf{Statement.}  The decomposition below identifies the mean
coordinate's warp cost with the level path's roughness,
$\dot{\bar u} = -\ddot\ell$.  Two residual contributions to
$\dot{\bar u}$ are set aside and accounted here: (a) the
\emph{common-mode disturbance response} --- every receiver's fast
loop responds to its own measured disturbance, and the fleet-average
of those responses (shared wander, shared measurement noise,
co-occurring pauses; \cref{as:cross}) is mean warp motion belonging
to neither the difference modes nor the level path; (b) \emph{drift-rate
variation}, $\dot{\bar w} \ne 0$.\\
\textbf{Role.}  Keeps $J_{\mathrm{diff}} + J_{\mathrm{mean}}$
additive in \cref{thm:structure}.\\
\textbf{Defect.}  (a) is computable from the transfer functions of
\cref{prop:transfer} evaluated on the \emph{common} component's
spectrum: the omitted cost is
$r\int |s\,U/Y(j\omega)|^2\, S_{\mathrm{common}}(\omega)\,
\dd\omega/2\pi$ for the relevant backend input, the same integral
that prices per-receiver roughness, just fed the common spectrum ---
the ``warp/sync leakage'' entries of \cref{tab:ledger} name exactly
this mechanism.  Under the latency measurement on disciplined hosts
the common driver is small; under the margin measurement co-occurring
pauses make it the dominant residual.  (b) is bounded by clock
physics: quartz frequency drift is parts in $10^{10}$ per second
(thermal); taking a two-orders-of-magnitude margin,
$|\dot w_i| \le 10^{-8}\,\mathrm{s}^{-1}$, the contribution
$r\,\E[\dot{\bar w}^2] \le 10^{-16}\,r$ is negligible at any sane
$r$.\\
\textbf{Monitor.}  Cross-receiver coherence of the loop innovations
(estimates $S_{\mathrm{common}}$); the fleet-average warp-rate
spectrum against the level layer's commanded $\ddot\ell$.
\end{simplification}

\subsection{The decomposition theorem}\label{sec:decomposition-theorem}

\begin{definition}[Layered policies]\label{def:layered}
A policy is \emph{layered} when it consists of: a level path
$\ell(t)$, adapted to the sender's information and broadcast to the
receivers; and per-receiver fast laws that act causally on the
receiver's own measured error relative to the broadcast level.  Write
$\mathcal{P}_{\mathrm{L}}$ for this class.  A layered policy is
specified by the pair $(\pi_{\mathrm{fast}}, \ell(\cdot))$.
\end{definition}

\begin{theorem}[Structure of the optimal controller]\label{thm:structure}
Under S1--S13:
\begin{enumerate}[label=(\alph*)]
  \item[(0)] \textbf{Decomposition.}  For every layered policy
    $\pi = (\pi_{\mathrm{fast}}, \ell(\cdot)) \in
    \mathcal{P}_{\mathrm{L}}$,
    \[
      J(\pi)
      \;=\;
      \underbrace{J_{\mathrm{diff}}\bigl(\pi_{\mathrm{fast}}\bigr)}_{
        \text{difference modes}}
      \;+\;
      \underbrace{J_{\mathrm{mean}}\bigl(\ell(\cdot)\bigr)}_{
        \text{level path}}
      \;+\; \Delta_{\mathrm{S}}(\pi),
    \]
    where the evaluation error $\Delta_{\mathrm{S}}(\pi)$, in
    loss-rate units ($\mathsf{u}/\mathrm{s}$), is bounded by the sum
    of the S1--S13 defects, each converted to an absolute loss-rate contribution
    (relative defects multiply the term they qualify; the register of
    \cref{sec:register} lists both forms).
\end{enumerate}
Within $\mathcal{P}_{\mathrm{L}}$:
\begin{enumerate}[label=(\roman*)]
  \item $J_{\mathrm{diff}}$ depends only on the fast-layer feedback
    and is minimized, independently of the level path and within the
    policy class of S10, by the law $v = \kappa^2 e - \sqrt2\kappa u$
    of \cref{prop:riccati} applied in each difference coordinate
    (equivalently: by every receiver running that law against the
    common broadcast level), with $\kappa = (q/r)^{1/4}$.  Under exact
    error observation the class restriction is vacuous and the law is
    optimal among all causal laws
    (\cref{prop:riccati,rem:gaussian-role}).
  \item $J_{\mathrm{mean}}$ depends only on the level path:
    \begin{equation}\label{eq:mean-objective}
      J_{\mathrm{mean}}
      = \E\left[\,\alpha\,\ell
        \;+\; \frac{c}{N}\sum_{i=1}^N R_i\bigl(\ell - \hat B_i\bigr)
        \;+\; r\,\Bigl(\frac{\dd^2\ell}{\dd t^2}\Bigr)^{\!2}\,\right]
    \end{equation}
    under the stationary law of the belief process, and is minimized
    \emph{within the hold-and-glide subclass} --- hold $\ell$ at the
    argmin of the static program (\cref{sec:level-program});
    rest-to-rest glides per \cref{sec:glide} when the argmin moves ---
    up to the S6 and S7 defects.  Anticipatory level policies, which
    act on forecasts of the belief rather than its current value (the
    direction endorsed by \cref{prop:ascent}), are not derived here.
\end{enumerate}
\end{theorem}

\begin{proof}
Assembly of proven pieces.  By \cref{prop:loss-coords} the pointwise
loss splits into mean terms, difference terms, and the indicator.  By
\cref{prop:dynamics-decouple} the coordinates have decoupled dynamics,
so a policy may be specified per coordinate; conversely any policy's
loss evaluates coordinate-wise on the split terms.  The difference
part is the stationary expectation (S1) of the quadratic cost
\cref{eq:lqr-setup} per coordinate, minimized by \cref{prop:riccati}
within the policy class of S10; the minimizing gains do not depend on
the disturbance covariances, so the same law is optimal in every
difference coordinate regardless of cross-receiver correlation.  The
indicator, after S2 (additive evaluation), S3--S4 (no fast-policy
dependence), S5 (compounding closed over the gap law) and S6 (levels
static on episode timescales), equals
$\tfrac{c}{N}\sum_i R_i(\mu_i)$ with $R_i$ from
\cref{prop:risk-evaluation} --- a function of the level path alone.
The remaining mean-coordinate quadratic cost is
$r\,\E[\dot{\bar u}^2]$; on level glides
$\dot{\bar u} = -\ddot\ell$ up to the two residuals of S13
(drift-rate variation and the common-mode disturbance response), so
its stationary cost is the third term of \cref{eq:mean-objective}
plus S13's registered contribution (S6 separates the timescales so no
cross-frequency terms survive).  The per-receiver level tilts are set
to zero by S11 (\cref{sec:tilt}); every dropped term is one of the
S-defects, collected in $\Delta_{\mathrm{S}}$.
\end{proof}

\begin{remark}[The optimality gap of the layered class --- an
assembled claim, not a theorem]\label{rem:gap-claim}
We also claim: the best layered policy is within the same defect
budget of the best causal policy.  The reasoning: the mechanisms an
unrestricted policy could exploit beyond the layered class are the
ones the simplifications discard, and each S-item bounds its
mechanism's worth (S3: in-event control; S10: filtering and nonlinear
response to the measurement; S11: per-receiver levels; S13:
coordinated common-mode response; S7: belief-coupled control).  This
is an \emph{assembled} bound --- each piece rigorous under its own
stated conditions, summed --- and the completeness of the mechanism
list is an engineering judgment, not a proven exhaustion of policy
space.  It sits outside the theorem for exactly that reason.
\end{remark}


\subsection{The level program}\label{sec:level-program}

\begin{proposition}[Operating point]\label{prop:foc}
The static part of \cref{eq:mean-objective} is minimized at the level
$\ell^*$ solving
\begin{equation}\label{eq:foc}
  \frac{c}{N}\sum_{i=1}^{N} g_i\bigl(\ell^* - \hat B_i\bigr)
  \;=\; \alpha
  \qquad\Longleftrightarrow\qquad
  \frac1N\sum_{i=1}^{N} g_i(\mu_i^*) \;=\; \frac{\alpha}{c}\,,
\end{equation}
with $g_i$ the occupation density (\cref{def:occupation}); under S8
and the interiority condition \cref{eq:interior-condition} the
minimizer exists, is unique, and satisfies \cref{eq:foc}
(\cref{prop:risk-shape}).  Only the ratio $\alpha/c$ enters: the operating point balances the summed
occupation density of the deviation at the margin against the price
ratio of latency to underrun.  Dimensional check: $[\alpha/c] =
\mathrm{s}^{-1} = [g_i]$ (\cref{rem:units}).
\end{proposition}

\begin{proof}
Differentiate $\alpha\ell + \tfrac{c}{N}\sum_i R_i(\ell - \hat B_i)$
in $\ell$ and use $R_i' = -g_i$; convexity (\cref{prop:risk-shape})
makes the stationary point the minimum.
\end{proof}

\begin{remark}[The target has emerged]\label{rem:emergence}
No reference latency appears anywhere in
\cref{sec:model,sec:process,sec:loss}.  The quantity $\ell^*$ of
\cref{eq:foc} is produced by the optimization: it is a function of the
estimated baseline levels and risk curves --- that is, of the
controller's belief state --- and it moves when the belief moves.  The
``latency target'' of the implementation is thus a \emph{dynamic state
variable of the controller}, with dynamics inherited from the
estimators through the argmin map, rate-limited by \cref{sec:glide}.
This is the promised answer to the question the problem statement left
open: a reference \emph{may} arise in the solution, and it does ---
as internal state, not as input.
\end{remark}

\subsection{Per-receiver tilts}\label{sec:tilt}

Allowing each receiver its own level $\ell_i$ turns the static program
into
$\min_{\ell_1,\dots,\ell_N} \alpha\bar\ell
+ \tfrac{q}{N}\sum_i(\ell_i-\bar\ell)^2
+ \tfrac{c}{N}\sum_i R_i(\ell_i - \hat B_i)$.

\begin{proposition}[Optimal tilts]\label{prop:tilt}
The stationary conditions are, for each $i$,
\[
  \frac{\alpha}{N} + \frac{2q}{N}\,(\ell_i - \bar\ell)
  \;=\; \frac{c}{N}\, g_i(\mu_i),
\]
whose average over $i$ recovers \cref{eq:foc}, and whose centered part
gives
\begin{equation}\label{eq:tilt}
  \ell_i - \bar\ell
  \;=\;
  \frac{c\,\bigl[g_i(\mu_i) - \tfrac1N\sum_j g_j(\mu_j)\bigr]}{2q}\,.
\end{equation}
Each receiver is displaced by its \emph{excess} risk pressure relative
to the fleet average, divided by twice the synchronization stiffness.
\end{proposition}

\begin{proof}
Differentiate in $\ell_i$:
$\partial_{\ell_i}\bigl[\tfrac{q}{N}\sum_j(\ell_j-\bar\ell)^2\bigr]
= \tfrac{2q}{N}(\ell_i-\bar\ell)$ (the cross terms cancel because
$\sum_j(\ell_j-\bar\ell)$ is invariant to $\ell_i$ shifts through
$\bar\ell$; expanding,
$\partial_{\ell_i}\sum_j(\ell_j-\bar\ell)^2
= 2(\ell_i-\bar\ell)(1-\tfrac1N)
- \tfrac2N\sum_{j\ne i}(\ell_j-\bar\ell)
= 2(\ell_i-\bar\ell)$).  The latency term contributes $\alpha/N$ and
the risk term $-\tfrac{c}{N}g_i$.  Averaging the conditions eliminates
the centered term; subtracting the average isolates \cref{eq:tilt}.
\end{proof}

\setcounter{simplification}{10}%
\begin{simplification}[S11: tilts set to zero]\label{simp:tilt}
\textbf{Statement.}  The broadcast level is common:
$\ell_i \equiv \ell^*$.\\
\textbf{Role.}  One number to distribute; receivers need no
per-receiver risk data.\\
\textbf{Defect.}  By \cref{eq:tilt} the forfeited displacement is of
order $c\,\Delta g/(2q)$ with $\Delta g$ the cross-receiver spread of
risk pressure; at the operating point $c\,g_i \sim \alpha$, so the
tilt scale is $\alpha/(2q)$ --- with any synchronization weight strong
enough to hold sub-millisecond skew, microseconds.  The forfeited loss
is quadratic in the tilt: with every tilt at its scale,
$\tfrac{q}{N}\sum_i(\text{tilt}_i)^2 \sim \alpha^2/(4q)$, negligible
against every other defect.\\
\textbf{Monitor.}  The spread of estimated $c\,g_i(\mu_i)$ across
receivers, reported next to $\alpha$.
\end{simplification}

\subsection{Moving the level: glides and asymmetry}\label{sec:glide}

The argmin $\ell^*(\hat\theta)$ moves when the belief moves.  Moving
the fleet level at rate $\gamma$ means every receiver warps by
$\gamma$: the glide is executed by the existing fast loops tracking the
moving broadcast and is bounded by $|\gamma| \le u_{\max}$.  Under the
loss, a \emph{steady} glide is free --- the warp term charges
$\dot{\bar u}^2 = \ddot\ell^{\,2}$, and a constant-rate glide has
$\ddot\ell = 0$.  What a level move costs is therefore: the
time-integral of the latency overhang (descents) or of the excess risk
(ascents), plus the roughness of \emph{changing} the glide rate.  The
descent problem has a closed-form solution.

\begin{proposition}[Optimal descent]\label{prop:descent}
Suppose the belief lowers the argmin by $\Delta$: the current level
sits $\Delta$ above the target.  Write the level relative to the
target, $\phi(t) \coloneqq \ell(t) - \ell^*$, and consider rest-to-rest
descents: $\phi(0)=\Delta$, $\dot\phi(0)=0$, $\phi(T)=0$,
$\dot\phi(T)=0$, with the duration $T$ free, subject to no undershoot
($\phi \ge 0$; below the target the risk term rules, so overshoot
credit is not available).  The cost of a descent path is
\begin{equation}\label{eq:descent-cost-functional}
  C[\phi] \;=\; \int_0^T \Bigl[\,\alpha\,\phi(t)
    + r\,\ddot\phi(t)^2\,\Bigr]\dd t
\end{equation}
(excess latency plus warp roughness; risk never enters because
$\phi \ge 0$ keeps the margin at or above its target value throughout,
and $R_i$ is nonincreasing).  The optimal descent is the quartic arc
\begin{equation}\label{eq:descent-arc}
  \phi(t) \;=\; \Delta\,\Bigl(1 - \tfrac{t}{T^*}\Bigr)^{\!3}
    \Bigl(1 + \tfrac{3t}{T^*}\Bigr),
  \qquad
  \boxed{\;T^* \;=\; 2\sqrt3\,\Bigl(\frac{r\,\Delta}{\alpha}
    \Bigr)^{\!1/4}\;}
\end{equation}
with monotone level, peak descent rate and total cost
\begin{equation}\label{eq:descent-numbers}
  \gamma_{\mathrm{pk}}
  \;=\; \frac{8\sqrt3}{27}\,\Delta^{3/4}
    \Bigl(\frac{\alpha}{r}\Bigr)^{\!1/4}
  \quad\text{at } t = T^*/3,
  \qquad
  C^* \;=\; \frac{16\sqrt3}{15}\,
    \Delta^{5/4}\,\alpha^{3/4}\,r^{1/4},
\end{equation}
the cost splitting $3{:}1$ between latency overhang and warp
roughness.  This holds when $\gamma_{\mathrm{pk}} \le u_{\max}$, i.e.,
for $\Delta \le \Delta_{\mathrm{sat}} \coloneqq
\bigl(27 u_{\max}/(8\sqrt3)\bigr)^{4/3} (r/\alpha)^{1/3}$; the
saturated case follows the proposition.
\end{proposition}

\begin{proof}
\emph{Euler--Lagrange.}  The integrand
$F(\phi, \ddot\phi) = \alpha\phi + r\ddot\phi^2$ has
$\partial F/\partial\phi = \alpha$,
$\partial F/\partial\dot\phi = 0$,
$\partial F/\partial\ddot\phi = 2r\ddot\phi$.  For a second-order
Lagrangian with $\phi$ and $\dot\phi$ fixed at both ends, the
Euler--Lagrange equation is
$\partial_\phi F - \tfrac{\dd}{\dd t}\partial_{\dot\phi}F
+ \tfrac{\dd^2}{\dd t^2}\partial_{\ddot\phi}F = 0$, i.e.,
$\alpha + 2r\,\phi^{(4)} = 0$: constant fourth derivative
$\phi^{(4)} = -a$ with $a \coloneqq \alpha/(2r)$.

\emph{The quartic family.}  Integrating and imposing
$\phi(0)=\Delta$, $\dot\phi(0)=0$, $\phi(T)=0$, $\dot\phi(T)=0$:
\[
  \phi(t) = -\frac{a t^4}{24} + \frac{c_3 t^3}{6}
    + \frac{c_2 t^2}{2} + \Delta,
  \qquad
  c_3 = \frac{aT}{2} + \frac{12\Delta}{T^3},
  \qquad
  c_2 = -\frac{aT^2}{12} - \frac{6\Delta}{T^2}.
\]
For each fixed $T$ the functional $C[\phi]$ is strictly convex in
$\phi$ (a positive quadratic form in $\ddot\phi$ plus a linear
term), so the Euler--Lagrange solution is the unique global minimizer
among all paths with these boundary values --- comparing family
members across $T$ therefore compares true per-$T$ optima.
Substituting and integrating termwise gives the cost of the family
member with duration $T$:
\begin{equation}\label{eq:descent-cost-T}
  C(T) \;=\; \frac{12\,\Delta^2 r}{T^3}
    \;+\; \frac{\alpha\,\Delta\,T}{2}
    \;-\; \frac{\alpha^2\,T^5}{2880\,r}.
\end{equation}

\emph{No interior minimum; the boundary of admissibility.}
Differentiate and substitute $x = T^4$:
$C'(T) = \tfrac{\alpha\Delta}{2} - \tfrac{36\Delta^2 r}{x}
- \tfrac{\alpha^2 x}{576 r}$.
As a function of $x > 0$ this is concave with maximum at
$x = 144\,\Delta r/\alpha$ (set the derivative
$36\Delta^2 r/x^2 - \alpha^2/(576r)$ to zero), where its value is
$\tfrac{\alpha\Delta}{2} - \tfrac{\alpha\Delta}{4}
- \tfrac{\alpha\Delta}{4} = 0$.  Hence $C'(T) \le 0$ for all $T$,
with equality exactly at $T^4 = 144\,\Delta r/\alpha$, i.e., at
$T = T^*$ of \cref{eq:descent-arc}: the cost \emph{decreases} in $T$
everywhere else.  The minimizer must therefore sit at the largest
admissible $T$.

\emph{Admissibility ends at $T^*$.}  Factor the velocity:
$\dot\phi(t) = -\tfrac{a}{6}\,t\,(t - T)(t - t_1)$ with
$t_1 = \tfrac{T}{2} + \tfrac{36\Delta}{a T^3}$ (the cubic
$\dot\phi$ has roots at $0$ and $T$ by construction; the product and
sum of the quadratic factor's roots give $t_1$).  Now
$t_1 \ge T \iff 72\Delta \ge aT^4 \iff T \le T^*$.  For
$T \le T^*$: on $(0, T)$ the factors give
$t > 0$, $(t{-}T) < 0$, $(t{-}t_1) < 0$, so $\dot\phi < 0$: the
descent is monotone and $\phi \ge 0$ holds.  For $T > T^*$:
$t_1 < T$, so $\dot\phi > 0$ on $(t_1, T)$; with $\phi(T) = 0$ this
means $\phi < 0$ just before $T$ --- the interior path undershoots
and is inadmissible.  One admissible alternative remains at long
horizons: touch $\phi = 0$ with $\dot\phi = 0$ at some $T' < T$ and
rest there.  Resting at the target costs nothing, so such a path is
exactly a descent of duration $T'$; admissibility forces
$T' \le T^*$, and $C$ decreasing on $(0, T^*]$ gives
$C(T') \ge C(T^*)$.  Hence the constrained optimum is exactly
$T = T^*$.

\emph{Values at $T^*$.}  At $T = T^*$, $t_1 = T^*$ and
$\dot\phi = -\tfrac{a}{6}\,t\,(t - T^*)^2$: the arrival is
tangential.  Integrating gives the closed form
\cref{eq:descent-arc} (check: $\phi(0)=\Delta$, and expanding
$\Delta(1 - t/T^*)^3(1 + 3t/T^*)$ reproduces the quartic with
$\phi^{(4)} = -24\Delta\cdot 3/T^{*4} = -72\Delta/T^{*4} = -a$ by
$T^{*4} = 144\Delta r/\alpha$ and $a = \alpha/2r$~$\checkmark$).
The peak rate: $\tfrac{\dd}{\dd t}[t(t-T^*)^2] = (t-T^*)(3t-T^*)$
vanishes at $t = T^*/3$, where
$\dot\phi = -\tfrac{a}{6}\cdot\tfrac{T^*}{3}\cdot
\tfrac{4T^{*2}}{9} = -\tfrac{2aT^{*3}}{81}$; substituting $a$ and
$T^*$ gives $\gamma_{\mathrm{pk}}$ of \cref{eq:descent-numbers}.
The cost: substituting $T^*$ into \cref{eq:descent-cost-T} gives
$C^* = \tfrac{16\sqrt3}{15}\Delta^{5/4}\alpha^{3/4}r^{1/4}$, of
which the latency part is $\tfrac{4\sqrt3}{5}(\cdot)$ and the
roughness part $\tfrac{4\sqrt3}{15}(\cdot)$: ratio $3{:}1$.
\end{proof}

\begin{remark}[Saturated descents]\label{rem:descent-saturated}
For $\Delta > \Delta_{\mathrm{sat}}$ the arc's peak rate would exceed
the clamp.  The constrained-optimal profile then cruises at
$u_{\max}$ with entry and exit segments; we do not derive the exact
constrained path (a bang--singular problem) and instead use: cruise at
$u_{\max}$, enter and exit along the unconstrained arcs of effective
size $\Delta_{\mathrm{sat}}$.  Its cost is
$\alpha\Delta^2/(2u_{\max})\,(1 + o(1))$ with the corner cost bounded
by $2\,C^*(\Delta_{\mathrm{sat}})$ --- an explicitly declared
approximation, conservative in the loss.
\end{remark}

\begin{proposition}[Ascent is always at the actuator bound]
\label{prop:ascent}
Suppose the belief raises the argmin by $\Delta$ (the current level
sits $\Delta$ too low), and call the deficit \emph{material} when its
integrated excess risk at any sub-saturated ascent exceeds the corner
cost $2\,C^*(\Delta_{\mathrm{sat}})$.  During the ascent the system
runs excess risk
$\tfrac{c}{N}\sum_i\bigl[R_i(\mu_i(t)) - R_i(\mu_i^*)\bigr] \ge 0$,
which decreases in the ascent rate for every $t$.  Under the loss, the
warp term charges the ascent nothing while the rate is held constant
($\ddot\ell = 0$); it charges only the entry and exit of the glide, a
bounded one-time roughness cost that does not grow with $\Delta$.
Hence, for a material deficit, the optimal ascent runs at
$\gamma = u_{\max}$ until the level is restored, up to the bounded
corner cost.  For small deficits near the optimum --- where the risk
gradient is near zero --- the soft mirror of the descent arc
(\cref{prop:descent} with the roles of the endpoint costs exchanged)
is the better move; saturating the clamp is only warranted when the
risk at stake dominates the corners.  The full ascent takes $\Delta/u_{\max}$ --- $2$\,s per
$10$\,ms of level --- and this delay is irreducible: no admissible
policy closes a level deficit faster.  Consequently the level layer
must be \emph{predictive} where possible: a regime change detected
early (a leading indicator of rising $\lambda_i$ or $\sigma_i$) buys
margin at the only speed the actuator allows.
\end{proposition}

\begin{proof}
Risk excess decreases pointwise in the rate because a faster ascent
reaches every intermediate level sooner and $R_i$ is nonincreasing in
the level.  The corner cost is bounded exactly as in
\cref{rem:descent-saturated}: entry and exit along unconstrained arcs
cost at most $2\,C^*(\Delta_{\mathrm{sat}})$, independent of
$\Delta$.  The time bound is the actuator constraint
$|\dot\ell| = |\bar u - \text{drift}| \le u_{\max} + O(10^{-4})$ from
\cref{prop:latency-dynamics}.
\end{proof}

\begin{remark}[What asymmetry is and is not derived]
\label{rem:asymmetry}
Three asymmetries between raising and lowering the level follow from
the loss; one popular claim does not, and we state both sides.

\emph{Derived.}  (1) \emph{Stakes}: above the optimum the objective's
slope is $\alpha - \tfrac{c}{N}\sum_i g_i$, which is zero at the
optimum and \emph{bounded by} $\alpha$, approaching $\alpha$ only well
above it; below the optimum the slope reaches the occupation-density
scale, which can exceed $\alpha$ by orders of magnitude.  Sitting high
is cheap; sitting low is not.
(2) \emph{Rates}: an ascent with a material deficit always saturates
the clamp (\cref{prop:ascent}); a descent follows the unconstrained
arc of \cref{prop:descent} whenever its peak rate
$\tfrac{8\sqrt3}{27}\Delta^{3/4}(\alpha/r)^{1/4}$ is below the clamp
--- and this is slower for cheaper latency, $T^* \propto
\alpha^{-1/4}$.  The rate asymmetry is real but conditional: it comes
from the latency--roughness tradeoff, not from any penalty on steady
glides, and for large $\Delta$ both directions saturate.
(3) \emph{Decisions under uncertainty} (S7): delaying a descent by
one second costs at most $\alpha\Delta$; delaying an ascent costs the
integrated excess risk, which can be far larger.  With a noisy
argmin estimate the optimal trigger is therefore asymmetric: raise on
weak evidence, lower only on statistically confident evidence.

\emph{Not derived.}  The loss contains no penalty on a steady glide
and no rate limit other than the clamp; any policy that descends
slower than \cref{prop:descent} prescribes is conservatism imported
from outside this loss, not a consequence of it.
\end{remark}

\subsection{What remains of the general problem}\label{sec:residual}

Three couplings survive the decomposition, each with a bounded loss:
in-event control (S3's defect --- grows with warp authority);
estimation--control coupling (S7's defect --- shrinks with data); and
the distributed-information gap: the theorem's fast layer needs only
the broadcast level, which tolerates the message delay because the
level is slow (S6), but a formally optimal distributed policy for the
transient case (one receiver disturbed, the fleet informed late) is
not derived here.  Team-decision problems of this kind have no general
closed form (imported: the counterexample literature on distributed
linear-quadratic control); the timescale argument bounds what the gap
can cost, and \cref{sec:limits} records it as open.

%% =============================================================================
\section{The Resulting Controller}\label{sec:controller}

\Cref{thm:structure} licenses a controller with four layers.  Nothing in
the fast layer is new relative to the deployed system; the level layer
and its broadcast are the new mechanism.

\subsection{Layer 1: estimators (the belief state)}
\label{sec:layer-estimators}

Per receiver, maintained by the telemetry already in production:
\begin{itemize}
  \item Baseline $\hat B_i$: controller-matched exponential mean of the
    delay level.  Its time constant matches the fast loop so that
    ``deviation'' means exactly the part of the delay the loop does not
    absorb.
  \item Deviation scale $\hat\sigma_i$ and the deviation histogram:
    estimates the bulk exceedance term of \cref{eq:risk-evaluation};
    frozen during detected episodes so the bulk estimate is not
    contaminated by the tail it gates.
  \item Episode meter: rate $\hat\lambda_i$, height histogram
    $\hat F_i$, gap histogram $\hat G_i$, per-episode recovery
    $\hat\rho$.  Every episode's height is measured whether or not it
    underruns, so the risk curve is identified at all levels without
    deliberately lowering the margin (S7).
  \item Drain validator: windowed margin minimum --- the model-free
    check of the composed $\hat R_i$ (monitor of S2, S5).
\end{itemize}

\subsection{Layer 2: the level program}\label{sec:layer-level}

At a slow cadence (seconds), the sender-side computes:
\begin{enumerate}
  \item Compose $\hat R_i(\mu)$ from \cref{eq:risk-evaluation} using
    layer-1 estimates; overlay the raw deviation-histogram exceedance
    where it has support (S2 monitor).
  \item Solve \cref{eq:foc} for $\ell^*$ by bracketing (S8): find the
    level where the fleet-mean occupation density crosses $\alpha/c$.
  \item Rate-limit the output: descend along the quartic arc of
    \cref{prop:descent} (duration $2\sqrt3\,(r\Delta/\alpha)^{1/4}$,
    peak rate capped at $u_{\max}$ per
    \cref{rem:descent-saturated}); ascend at $u_{\max}$
    (\cref{prop:ascent}), triggered eagerly on regime-change
    detection (S1 monitor); gate descents on estimation confidence
    (S7, \cref{rem:asymmetry}).
\end{enumerate}
The single number $\ell^*(t)$, plus per-receiver baseline offsets where
backends differ, is broadcast to receivers.  This is the only new
protocol surface the theory requires.

\subsection{Layer 3: fast loops (unchanged)}\label{sec:layer-fast}

Each receiver tracks the broadcast level with the law of
\cref{prop:riccati} plus integral action \cref{eq:pid}:
\[
  v = \kappa^2 e + K_0 \xi - \sqrt2\,\kappa\, u,
  \qquad \kappa = (q/r)^{1/4},\ \zeta = 1/\sqrt2,
\]
with anti-windup at the warp clamp.  This is the deployed second-order
controller; \cref{rem:deployed} maps the constants.  By
\cref{thm:structure}(i) it is the optimum for its layer --- no redesign
is licensed by the joint problem.

\subsection{Layer 4: event hygiene}\label{sec:layer-hygiene}

During detected episodes: freeze the bulk estimators (layer 1), do not
attempt heroic warp responses (S3: the authority is not there), record
the episode's height, gap, and recovery for the risk curve.  After the
episode: let the redelivery burst restore the margin where the
mechanism provides it; where recovery is warp-limited, the elevated
compound risk during $\rho$ is already priced by
\cref{eq:risk-evaluation}.

\subsection{Self-monitoring}\label{sec:self-monitoring}

Each simplification carries a live validity ratio
(\cref{tab:register}).  The controller can therefore report not only
its decisions but the current size of the modeling error behind them.
Two monitors deserve emphasis: the S3 ratio
$u_{\max}\mu\,g(\mu)/R(\mu)$, because the entire decomposition stands
on it, and the drain validator, because it checks the composed risk
curve against reality without using the model.

%% =============================================================================
\section{Register of Simplifications}\label{sec:register}

\begin{table}[h!]
\centering
\caption{All simplifications, their defects, and their monitors.
Reference numbers use $u_{\max}=5\times10^{-3}$, margins
$5$--$100$\,ms, $\sigma_i$ up to $1$\,ms wired, heights up to $1$\,s.}
\label{tab:register}
\small
\begin{tabular}{@{}p{0.34\textwidth}p{0.34\textwidth}p{0.24\textwidth}@{}}
\toprule
Simplification & Defect (definition and size) & Monitor \\
\midrule
S1 stationary ergodic regime
& loss during regime adaptation: switch rate $\times$ adaptation time
  $\times$ excess loss
& short- vs long-window estimator disagreement \\
\addlinespace
S2 additive risk evaluation
& threshold smearing $O(\sigma_i/\text{height spread})$; duty cycle
  $O(\lambda\E H)$; zero-drain ramp (burst drain time)
& composed $\hat R$ vs raw deviation histogram \\
\addlinespace
S3 fast layer cannot move the indicator
& absolute $\le \alpha\,u_{\max}\mu$ ($\le u_{\max}$ of the latency
  term); relative to the risk term: $u_{\max}\mu/x_0$, up to
  $\sim\!5\%$ for the bulk part
& ratio $u_{\max}\mu\,g(\mu)/R(\mu)$; anchor $\alpha u_{\max}\mu$ \\
\addlinespace
S4 frozen margin during episode
& subsumed in S3
& same \\
\addlinespace
S5 one-step compounding
& chain terms $O((\lambda^{\mathrm{loc}}\rho)^2)$; gap--height
  covariance
& episodes within $\rho$ of predecessor; drain validator \\
\addlinespace
S6 two-timescale separation
& glide-concurrent episodes: margin error up to $\gamma H \le 5$\,ms
& episodes stamped with glide rate \\
\addlinespace
S7 certainty equivalence
& $\tfrac12 J''(\mu^*)\Var[\hat\mu^*]$; shrinks as $1/n$ with event
  count, inflated by clustering
& confidence band on $\hat\mu^*$ \\
\addlinespace
S8 smooth, convex risk curve
& subgradient bracket at atoms; compounding can break convexity in
  the clustered regime --- program falls back to a 1-D scan, losing
  only the first-order characterization
& $\hat g$ near operating point; second difference of scanned
  $\hat J$ \\
\addlinespace
S9 component independence (2nd-order accounting)
& dropped cross-covariances; affects evaluations, not gains, not
  $\hat R$
& baseline-drift vs event cross-correlation \\
\addlinespace
S10 full-information law on the raw measurement
& backend-dependent: measurement-noise passthrough
  $\int|\kappa^2/\mathcal{D}|^2 S_n$ (latency meas.); forfeited
  nonlinear gain vs pause leakage (margin meas.)
& innovation histogram and spectrum; $S_n$ per host \\
\addlinespace
S11 zero tilts
& forfeited loss $\sim \alpha^2/4q$; tilt scale $\alpha/2q$
  (microseconds)
& spread of $c\,g_i(\mu_i)$ vs $\alpha$ \\
\addlinespace
S12 white-noise surrogate in the fast-loop design
& evaluation error
  $\int|H|^2(S_{\mathrm{true}}-\sigma_w^2)\,\dd\omega/2\pi$; the DC
  gain defect is repaired by integral action
& innovation-spectrum flatness within $\kappa$ \\
\addlinespace
S13 common-mode loop response priced separately
& omitted mean-warp roughness: the per-receiver roughness integral
  fed the common spectrum; drift-rate part $\le 10^{-16}r$
& cross-receiver coherence of loop innovations; fleet-average
  warp-rate spectrum vs commanded $\ddot\ell$ \\
\bottomrule
\end{tabular}
\end{table}

%% =============================================================================
\section{Limits of the Result}\label{sec:limits}

Stated as openly as the simplifications.

\begin{enumerate}[label=(L\arabic*)]
  \item \textbf{Warp authority.}  The decomposition stands on S3.  If
    $u_{\max}$ were raised by an order of magnitude or more, in-event
    control would become worthwhile, the indicator would couple to the
    fast layer, and the correct formulation is a free-boundary dynamic
    program that this document does not solve.  The S3 monitor
    quantifies the approach to that regime.
  \item \textbf{Risk aversion.}  The loss charges mean silence time.
    A listener may care more about the worst minute than the average
    second.  A variance-sensitive or quantile objective re-imports the
    temporal law (clustering) into the \emph{objective} --- see
    \cref{rem:division} --- and the level program stops being a
    one-dimensional convex program in a single scalar.
  \item \textbf{Session supervision.}  Aborts and restarts
    (\cref{rem:no-restart}) are outside the loss.  In particular,
    implementations in which playout stalls instead of inserting
    silence violate \cref{conv:silence}: a stall converts an episode
    into a schedule slip whose recovery is warp-limited
    (\cref{as:recovery}(b)) and whose transient can trip supervision
    thresholds.  The convention--implementation gap is a real
    difference, not a modeling detail, and must be closed on the
    implementation side or modeled explicitly.
  \item \textbf{Distributed transients.}  The optimal reaction of the
    undisturbed majority to one disturbed receiver, under message
    delay, is not derived (\cref{sec:residual}).  The forfeited value
    is bounded by the disturbance rate times the reconvergence
    transient.
  \item \textbf{Lost content.}  Episodes whose content is lost rather
    than delivered late cause silence that no margin prevents; that
    channel belongs to forward error correction, not to latency
    control, and is additive to the loss studied here.
  \item \textbf{Adversarial or trending regimes.}  S1 covers slow
    regime switching with re-estimation.  Systematically trending
    networks (a slow decay toward failure) are tracked by the belief
    but never anticipated beyond the leading indicators of
    \cref{prop:ascent}.
\end{enumerate}

%% =============================================================================
\appendix

\section{Imported Results}\label{sec:app-imported}

\subsection{Campbell--Mecke formula}\label{sec:app-campbell}
For a stationary marked point process on $\R$ with intensity $\lambda$
and Palm mark distribution $F$, and any nonnegative measurable function
$f$ of the mark,
\[
  \E\Bigl[\sum_{k:\,T_k\in[0,T]} f(H_k)\Bigr] \;=\; \lambda\,T
  \int f\,\dd F .
\]
Reference: Daley and Vere-Jones, \emph{An Introduction to the Theory of
Point Processes}, Vol.~II, Ch.~13 (Palm theory); the ergodic form used
in \cref{prop:campbell} is the corresponding ratio limit theorem.  Note
what is \emph{not} required: no Poisson assumption, no independence
between points --- stationarity suffices, which is why clustering does
not move the mean in \cref{eq:campbell}.

\subsection{Rice's formula}\label{sec:app-rice}
For a stationary, almost surely continuously differentiable Gaussian
process $X$ with $\Var[X] = \sigma_0^2$ and $\Var[\dot X] = \sigma_1^2$,
the mean rate of upcrossings of level $z$ is
\[
  N^+(z) \;=\; \frac{\sigma_1}{2\pi\,\sigma_0}\,
  \exp\!\Bigl(-\frac{z^2}{2\sigma_0^2}\Bigr).
\]
Reference: Rice (1944); Leadbetter, Lindgren, Rootz\'en,
\emph{Extremes and Related Properties of Random Sequences and
Processes}, Ch.~7.  The same reference (Ch.~7, sojourn results) covers
the ergodic occupation--crossing ratio used alongside it: for a
stationary process with finitely many level crossings per unit time,
the mean duration of a visit above $z$ equals the stationary
exceedance probability divided by the upcrossing rate.  The sojourn
computation of \cref{prop:sojourn} combines exactly these two
imports.

\subsection{Gaussian tail independence}\label{sec:app-tail-independence}
\begin{lemma}\label{lem:tail-independence}
Let $(Z_1, Z_2)$ be bivariate standard normal with correlation
$\varsigma < 1$.  Then
$\lim_{z\to\infty} \Prob(Z_2 > z \mid Z_1 > z) = 0$.
\end{lemma}
\begin{proof}
For $\varsigma \le 0$, Slepian's comparison inequality (imported;
Slepian 1962) gives
$\Prob(Z_1 > z,\, Z_2 > z) \le \Prob(Z_1>z)\,\Prob(Z_2>z)$, so the
conditional probability is at most the unconditional one,
$\bar\Phi(z) \to 0$; assume therefore $0 < \varsigma < 1$.
Condition on $Z_1 = x$:
$Z_2 \sim \mathcal N(\varsigma x,\, 1-\varsigma^2)$, so
$\Prob(Z_2 > z \mid Z_1 = x)
= \bar\Phi\bigl((z-\varsigma x)/\sqrt{1-\varsigma^2}\bigr)$.
Split the conditioning event at $x = z(1+\epsilon)$ for a fixed
$\epsilon > 0$.  For $x \le z(1+\epsilon)$:
$(z - \varsigma x)/\sqrt{1-\varsigma^2}
\ge z\,(1 - \varsigma(1+\epsilon))/\sqrt{1-\varsigma^2} \to \infty$
when $\epsilon < (1-\varsigma)/\varsigma$, so the conditional
probability tends
to $0$ uniformly on this range.  For the range $x > z(1+\epsilon)$,
its share of the conditioning mass is
$\bar\Phi(z(1+\epsilon))/\bar\Phi(z) \to 0$ by the Mills ratio
$\bar\Phi(t) = \varphi(t)/t\,(1+O(t^{-2}))$, since
$\varphi(z(1+\epsilon))/\varphi(z) = e^{-z^2(2\epsilon+\epsilon^2)/2}
\to 0$.  Both contributions vanish.
\end{proof}
Reference: Sibuya (1960).  This is the mechanism behind
\cref{prop:sojourn}: a Gaussian-copula process cannot hold an extreme
level; a pause can and does.

\subsection{Other imports}\label{sec:app-other}
The verification theorem for average-cost linear-quadratic control
(\cref{prop:riccati}): Fleming and Soner, \emph{Controlled Markov
Processes and Viscosity Solutions}, or any standard treatment of ergodic
LQG.  The spectral representation of stationary processes
(\cref{prop:psd}): Wiener--Khinchin/Bochner, standard.  The
distributed-control counterexamples referenced in \cref{sec:residual}:
Witsenhausen (1968).

\section{Notation}\label{sec:notation}

\begin{table}[h!]
\centering
\small
\begin{tabular}{@{}llll@{}}
\toprule
Symbol & Meaning & Units & Defined \\
\midrule
$s$ & stream position & s (audio) & \cref{def:stream} \\
$\tau(s)$ & emission time of $s$ & s & \cref{def:stream} \\
$a_i(s)$ & arrival time at receiver $i$ & s & \cref{def:arrival} \\
$D_i(s)$ & delivery delay & s & \cref{def:arrival} \\
$p_i(s),\ s_i(t)$ & playout schedule and its inverse & s &
  \cref{def:schedule} \\
$u_i,\ u_{\max}$ & warp; warp clamp ($5\times10^{-3}$) & --- &
  \cref{def:schedule} \\
$v_i = \dot u_i$ & warp rate (control input) & $\mathrm{s}^{-1}$ &
  \cref{eq:lqr-setup} \\
$L_i,\ \Lbar,\ \Ltil{i}$ & latency; fleet mean; deviation & s &
  \cref{def:latency,def:coordinates} \\
$w_i$ & clock drift disturbance & --- & \cref{prop:latency-dynamics} \\
$m_i$ & margin & s & \cref{def:margin} \\
$B_i,\ \varepsilon_i,\ P_i$ & baseline, diffusion, pauses & s &
  \cref{as:decomposition} \\
$\sigma_i$ & diffusion scale & s & \cref{as:diffusion} \\
$(T_{i,k}, H_{i,k})$ & episode start, height & s & \cref{as:pause} \\
$\lambda_i,\ F_i,\ G_i$ & episode rate; height law; gap law &
  $\mathrm{s}^{-1}$, ---, --- & \cref{as:pause} \\
$\rho_{i,k},\ \chi_{i,k}$ & recovery time; consumed margin & s &
  \cref{def:recovery,eq:depression} \\
$x_i$ & deviation $\varepsilon_i + P_i$ & s & \cref{eq:margin-decomp} \\
$\alpha, q, c, r$ & loss weights &
  see \cref{rem:units} & \cref{def:loss} \\
$\kappa = (q/r)^{1/4}$ & fast-loop bandwidth & $\mathrm{s}^{-1}$ &
  \cref{prop:riccati} \\
$\ell,\ \ell^*$ & latency level; its optimum & s &
  \cref{sec:level-program} \\
$\mu_i$ & margin level (backend-conditional) & s &
  \cref{eq:margin-level} \\
$R_i(\mu),\ g_i(\mu)$ & risk curve; occupation density & ---,
  $\mathrm{s}^{-1}$ & \cref{eq:R-def,def:occupation} \\
$\gamma$ & level glide rate & --- & \cref{sec:glide} \\
$\mathcal{D}(s)$ (Laplace) & $s^2+\sqrt2\kappa s+\kappa^2$ & --- &
  \cref{prop:riccati} \\
$\sigma_w^2$ & white-noise intensity (wander $+$ measurement) &
  $\mathrm{s}^2/\mathrm{s}$ & \cref{sec:lqr} \\
$\Lambda(t,\omega)$ & pointwise loss rate & $\mathsf{u}/\mathrm{s}$ &
  \cref{eq:pointwise-loss} \\
$e, u, v$ & fast-loop error, warp, input (one coordinate) &
  s, ---, $\mathrm{s}^{-1}$ & \cref{eq:lqr-setup} \\
$\xi,\ K_0,\ T_I$ & integral state, gain, settling time &
  $\mathrm{s}^2$, $\mathrm{s}^{-3}$, s & \cref{eq:pid} \\
$\delta(g)$ & residual margin depression at gap $g$ & s &
  \cref{eq:depression} \\
$\Pi$ & centering matrix $I - \mathbf{1}\mathbf{1}^\top/N$ & --- &
  \cref{prop:modes} \\
$\bar\beta$ & baseline slope bound & --- & \cref{as:baseline} \\
$x_0, C$ & diffusion tail scale and constant & s, --- &
  \cref{as:diffusion} \\
$\Delta,\ \Delta_{\mathrm{sat}}$ & level change; saturation
  threshold & s & \cref{prop:descent} \\
$T^*,\ \gamma_{\mathrm{pk}},\ C^*$ & descent duration, peak rate,
  cost & s, ---, $\mathsf{u}$ & \cref{prop:descent} \\
$J,\ J_{\mathrm{diff}},\ J_{\mathrm{mean}}$ & loss and its
  decomposition parts & $\mathsf{u}/\mathrm{s}$ &
  \cref{def:loss,thm:structure} \\
$\tau_{\hat B}$ & baseline filter time constant & s &
  \cref{sec:indicator-exact} \\
$n_i,\ N_i$ & measurement error (time domain, transform) & s &
  \cref{def:meas-e2e,conv:transforms} \\
$\tilde D$ & transform of the deviation $x_i$ & s &
  \cref{conv:transforms} \\
$\hat\theta$ & belief state (estimator outputs) & --- &
  \cref{def:belief} \\
$e_i^{\mathrm{loop}}$ & fast-loop tracking error & s &
  \cref{eq:margin-decomp} \\
$b_i$ & baseline offsets in the level objective & s &
  \cref{prop:risk-shape} \\
$\lambda_i^{\mathrm{loc}}$ & local (in-cluster) episode rate &
  $\mathrm{s}^{-1}$ & \cref{simp:compound} \\
$\bar\Phi,\ \varphi$ & standard normal tail, density & --- & standard \\
$\varrho$ & average-cost constant & $\mathsf{u}/\mathrm{s}$ &
  \cref{prop:riccati} \\
$\varsigma$ & correlation (tail-independence lemma) & --- &
  \cref{lem:tail-independence} \\
$\Delta_{\mathrm{S}}(\pi)$ & evaluation error of the decomposition &
  $\mathsf{u}/\mathrm{s}$ & \cref{thm:structure} \\
\bottomrule
\end{tabular}
\end{table}

\end{document}
